\chapter{归纳}
\label{cha:induction}

在 \cref{cha:typetheory} 中，我们介绍了许多从已有类型构造新类型的方法。
除了（依值）函数类型和宇宙之外，所有这些规则都是\emph{归纳定义（inductive definition）}\index{definition!inductive}\index{inductive!definition}这一普遍概念的特殊情形。
在本章中，我们将更一般地研究归纳定义。

\section{归纳类型简介}
\label{sec:bool-nat}

\index{type!inductive|(}%
\index{generation!of a type, inductive|(}
直观上，一个\emph{归纳类型（inductive type）} $X$ 可以理解为由某个有限的\emph{构造子（constructors）}集合“自由生成”的类型；其中每个构造子都是一个以 $X$ 为陪域的函数（可以带有若干个参数）。
这也包括零个参数的函数，它们就是 $X$ 的元素。

在描述一个具体的归纳类型时，我们逐条列出其构造子。
例如，\cref{sec:type-booleans} 中的类型 \bool 是由以下两个构造子归纳生成的：

\begin{itemize}
\item $\bfalse:\bool$
\item $\btrue:\bool$
\end{itemize}

类似地，$\unit$ 由以下一个构造子归纳生成：

\begin{itemize}
\item $\ttt:\unit$
\end{itemize}

而 $\emptyt$ 则根本没有构造子。
构造子带参数的一个例子是余积 $A+B$，它由两个构造子生成：

\begin{itemize}
\item $\inl:A\to A+B$
\item $\inr:B\to A+B$。
\end{itemize}

构造子带多个参数的一个例子是 Cartesian 积 $A\times B$，它由一个构造子生成：

\begin{itemize}
\item $\pairr{\blank, \blank} : A\to B \to A\times B$。
\end{itemize}

至关重要的是，我们还允许归纳类型的构造子接受来自正在定义的归纳类型本身的参数。
例如，自然数类型 $\nat$ 具有构造子：

\begin{itemize}
\item $0:\nat$
\item $\suc:\nat\to\nat$。
\end{itemize}

\symlabel{lst}
\index{type!of lists}%
\indexsee{list}{type of lists}%
另一个有用的例子是由某类型 $A$ 的元素组成的有限列表类型 $\lst A$，它具有构造子：

\begin{itemize}\label{list_constructors}
\item $\nil:\lst A$
\item $\cons:A\to \lst A \to \lst A$。
\end{itemize}

\index{free!generation of an inductive type}%
直观上，我们应该将一个归纳类型理解为其构造子\emph{自由生成的}类型。
也就是说，归纳类型的元素恰好就是从无开始、反复应用构造子所能得到的一切。
（我们将在 \cref{sec:identity-systems,cha:hits} 看到，对于更一般的归纳定义，这个概念需要稍作修改，但目前这样理解已经足够。）
例如，对于 \bool，我们应当认为其元素只有 $\bfalse$ 和 $\btrue$。
类似地，对于 \nat，我们应当认为每个元素要么是 $0$，要么是对某个“先前构造的”自然数应用 $\suc$ 得到的。

\index{induction principle}%
然而，我们并非直接断言此类性质，而是通过\emph{归纳原理（induction principle）}，也称为\emph{（依值）消去规则（(dependent) elimination rule）}来表达它们。
我们在 \cref{cha:typetheory} 已经见过这些原理。
例如，\bool 的归纳原理是：
%

\begin{itemize}
\item 要证明关于 \bool 的\emph{所有}元素的命题 $E : \bool \to \type$，只需对 $\bfalse$ 和 $\btrue$ 分别证明即可，亦即给出证明 $e_0 : E(\bfalse)$ 和 $e_1 : E(\btrue)$。
\end{itemize}

此外，将所得的证明 $\ind\bool(E,e_0,e_1): \prd{b : \bool}E(b)$ 作用于构造子 $\bfalse$ 和 $\btrue$ 时，其行为符合预期；这一原理由\emph{计算规则（computation rules）}\index{computation rule!for type of booleans}表达：

\begin{itemize}
\item 我们有 $\ind\bool(E,e_0,e_1,\bfalse) \jdeq e_0$。
\item 我们有 $\ind\bool(E,e_0,e_1,\btrue) \jdeq e_1$。
\end{itemize}

\index{case analysis}%
因此，布尔类型 $\bool$ 的归纳原理允许我们进行\emph{分情形讨论（case analysis）}。
由于两个构造子都不带任何参数，这对布尔类型来说已经足够了。

\index{natural numbers}%
然而，对于自然数，仅靠分情形讨论通常是不够的：在对应于归纳步骤 $\suc(n)$ 的情形中，我们还需要假设待证的命题对 $n$ 已经成立。
这就得到了如下的归纳原理：

\begin{itemize}
\item 要证明关于\emph{所有}自然数的命题 $E : \nat \to \type$，只需对 $0$ 和 $\suc(n)$ 分别证明即可，其中在证明 $\suc(n)$ 时可以假设命题对 $n$ 成立，亦即我们构造 $e_z : E(0)$ 和 $e_s : \prd{n : \nat} E(n) \to E(\suc(n))$。
\end{itemize}

与布尔类型的情况类似，对于函数 $\ind\nat(E,e_z,e_s) : \prd{x:\nat} E(x)$，我们也有相应的计算规则：
\index{computation rule!for natural numbers}%

\begin{itemize}
\item $\ind\nat(E,e_z,e_s,0) \jdeq e_z$。
\item 对任意 $n : \nat$，有 $\ind\nat(E,e_z,e_s,\suc(n)) \jdeq e_s(n,\ind\nat(E,e_z,e_s,n))$。
\end{itemize}

因此，依值函数 $\ind\nat(E,e_z,e_s)$ 可以理解为根据参数 $x : \nat$ 递归定义的；这一定义通过函数 $e_z$ 和 $e_s$ 实现，我们称它们为\define{递归式（recurrences）}\indexdef{recurrence}。
当 $x$ 为零\index{zero}时，该函数直接返回 $e_z$。
当 $x$ 是另一个自然数 $n$ 的后继\index{successor}时，结果是这样得到的：取递归式 $e_s$，并将具体的前驱\index{predecessor} $n$ 以及递归调用的值 $\ind\nat(E,e_z,e_s,n)$ 代入其中。

上述所有例子的归纳原理都具有这种家族相似性。
在 \cref{sec:strictly-positive} 中，我们将讨论“归纳定义”的一般概念，以及如何为其推导出合适的\emph{归纳原理}；但首先，我们来探究归纳定义之间的若干共通之处。

\index{recursion principle}%
例如，我们曾在 \cref{cha:typetheory} 的每一个例子里指出，从归纳原理可以推导出一个\emph{递归原理}，其中值域是简单类型（而非类型族）。
归纳原理和递归原理可能看起来有些奇怪，因为它们似乎只给出了函数的\emph{存在性}，而没有刻画出其唯一性。
然而事实上，归纳原理本身已经足够强，足以证明其自身的\emph{唯一性原理}\index{uniqueness!principle, propositional!for functions on N@for functions on $\nat$}，如下述定理所示。

\begin{thm}\label{thm:nat-uniq}
设 $f,g : \prd{x:\nat} E(x)$ 为两个函数，它们在命题相等的意义下满足递归式
%
\begin{equation*}
  e_z : E(0)
  \qquad\text{and}\qquad
  e_s : \prd{n : \nat} E(n) \to E(\suc(n))
\end{equation*}
%
也就是说，满足
\begin{equation*}
  \id{f(0)}{e_z}
  \qquad\text{and}\qquad
  \id{g(0)}{e_z}
\end{equation*}
以及
\begin{gather*}
  \prd{n : \nat} \id{f(\suc(n))}{e_s(n, f(n))},\\
  \prd{n : \nat} \id{g(\suc(n))}{e_s(n, g(n))}.
\end{gather*}
那么 $f$ 和 $g$ 相等。
\end{thm}

\begin{proof}
我们在类型族 $D(x) \defeq \id{f(x)}{g(x)}$ 上使用归纳法。对于基始情形，有
\[ f(0) = e_z = g(0). \]
对于归纳情形，假设 $n : \nat$ 且 $f(n) = g(n)$。那么
\[ f(\suc(n)) = e_s(n, f(n)) = e_s(n, g(n)) = g(\suc(n)). \]
第一个和最后一个等式由 $f$ 和 $g$ 所满足的条件得出。中间的等式则由归纳假设以及函数应用保持相等性这一事实得出。这样我们就得到了 $f$ 与 $g$ 的逐点相等；再由函数外延性即可完成证明。
\end{proof}

注意，唯一性原理甚至适用于那些仅\emph{在命题相等意义下}满足递归式的函数，亦即一条道路。
当然，由归纳原理直接得到的那个特定函数是在判断相等性\index{judgmental equality}的意义下满足这些递归式的；我们将在 \cref{sec:htpy-inductive} 回到这一点。
另一方面，定理本身只断言了函数之间的一个命题相等（亦见 \cref{ex:same-recurrence-not-defeq}）。
从同伦的观点来看，很自然地会问这条道路是否\emph{相容}，即等式 $f=g$ 在高阶道路意义下是否唯一；在 \cref{sec:initial-alg} 我们将会看到事实的确如此。

当然，类似的函数唯一性定理通常也能为其它归纳类型表述并证明。
在下一节中，我们将展示这一唯一性性质如何与泛等性结合，从而推出像自然数这样的归纳类型完全由其引入规则、消去规则和计算规则所刻画。

\index{type!inductive|)}%
\index{generation!of a type, inductive|)}%

%%%%%%%%%%%%%%%%%%%%%%%%%%%%%%%%%%%%%%%%%%%%%%%%%%%%%%%%%%

\section{归纳类型的唯一性}
\label{sec:appetizer-univalence}

\index{uniqueness!of identity types}%
我们已将“那个”自然数定义为一个特定的类型 \nat，并配有特定的归纳构造子 $0$ 和 $\suc$。
然而，根据前一节描述的类型论中归纳定义的一般原理，我们完全可以用相同的方式定义\emph{另一个}类型。
\index{natural numbers!isomorphic definition of}
也就是说，假设我们令 $\natp$ 是由以下构造子生成的归纳类型：

\begin{itemize}
\item $\zerop:\natp$
\item $\sucp:\natp\to\natp$。
\end{itemize}

那么 $\natp$ 将具有与 $\nat$ 形式上相同的归纳原理和递归原理。
当要为所有这些“新”自然数证明一个命题 $E : \natp \to \type$ 时，只需给出证明 $e_z : E(\zerop)$ 和 \narrowequation{e_s : \prd{n : \natp} E(n) \to E(\sucp(n))。}
而函数 $\rec\natp(E,e_z,e_s) : \prd{n:\natp} E(n)$ 则具有如下计算规则：

\begin{itemize}
\item $\rec\natp(E,e_z,e_s,\zerop) \jdeq e_z$，
\item $\rec\natp(E,e_z,e_s,\sucp(n)) \jdeq e_s(n,\rec\natp(E,e_z,e_s,n))$ 对任意 $n : \natp$ 成立。
\end{itemize}

但是，$\nat$ 和 $\natp$ 之间有什么关系呢？

这并非仅仅是一个学术问题，因为许多其他地方也能找到``像''自然数的结构。
\index{type!of lists}%
例如，我们可以将自然数与单元素类型上的列表视为等同（这可以说是最古老的呈现方式——见于洞穴墙壁上\index{caves, walls of}），与非负整数、有理数和实数的某些子集等等视为等同。
而从编程\index{programming}的角度看，自然数的这种``一元''表示效率非常低，因此我们有时可能更倾向于使用二进制表示。
我们希望能够在所有这些``自然数''版本之间相互等同，以便将构造和结果从一个版本转移到另一个版本。

当然，如果自然数的两个版本满足完全相同的归纳原理，那么它们导出的结构也完全相同。
譬如，回顾在\cref{sec:inductive-types}中定义的函数 $\dbl$ 的例子。为我们的新自然数定义一个类似的函数，只需依样复制并加上撇号即可：
\[ \dblp \defeq \rec\natp(\natp, \; \zerop, \;  \lamu{n:\natp}{m:\natp} \sucp(\sucp(m))). \]
这看起来简单，却有一个明显的缺点：会导致大量重复。
不仅函数需要复制，所有的引理及其证明也得复制。
例如，像$\prd{n : \nat} \dbl(\suc(n))=\suc(\suc(\dbl(n)))$这样一个简单的结果，连同其归纳证明，也必须被``加上撇号''。

在传统数学中，人们直接宣称 $\nat$ 和 $\natp$ 显然``相同''，并且可以在需要时互相替换。
这通常没什么问题，但它把相当多的东西扫到了地毯下面，拉大了非形式化数学与其精确描述之间的差距。
在同伦类型论中，我们可以做得更好。

首先注意到我们有如下可定义的映射：

\begin{itemize}
\item $f \defeq \rec\nat(\natp, \; \zerop, \;  \lamu{n:\nat} \sucp)
       : \nat \to\natp$,
\item $g \defeq \rec\natp(\nat, \; 0, \;  \lamu{n:\natp} \suc)
       : \natp \to\nat$。
\end{itemize}

由于 $g$ 和 $f$ 的复合满足与 $\nat$ 上恒等函数相同的递归式，由\cref{thm:nat-uniq}可得$\prd{n : \nat} \id{g(f(n))}{n}$，而同一定理的``加撇''版本则给出$\prd{n : \natp} \id{f(g(n))}{n}$。
因此，$f$ 和 $g$ 互为拟逆，故有 $\eqv{\nat}{\natp}$。
现在，我们可以沿着这个等价，将 $\nat$ 上的函数直接转移到 $\natp$ 上（反之亦然），例如：
\[ \dblp \defeq \lamu{n:\natp} f(\dbl(g(n))). \]
证明这个版本的 $\dblp$ 与早先定义的相等，是一个简单的练习。

当然，这没什么好惊讶的；数学家正是通过这样的同构来设想如何将 $\nat$ 与 $\natp$``等同''起来的。
\index{isomorphism!transfer across}%
然而，通过同构进行``转移''的机制，取决于被转移的对象；它并不总是像用 $f$ 和 $g$ 对单个函数做前复合与后复合那么简单。
考虑，例如，这样一个简单的引理：
\[\prd{n : \natp} \dblp(\sucp(n))=\sucp(\sucp(\dblp(n))).\]
插入正确的 $f$ 和 $g$，比起直接对 $n:\natp$ 重新进行归纳证明，也好不了太多。

\index{univalence axiom}%
正是在此处，泛等公理发挥了作用：由于 $\eqv{\nat}{\natp}$，我们还有 $\id[\type]{\nat}{\natp}$，即 $\nat$ 和 $\natp$ 作为类型是\emph{相等}的。
现在，关于相等类型的归纳原理保证了任何涉及 $\nat$ 的构造或证明，都能以同样的方式自动转移到 $\natp$ 上。
我们只需将该函数或定理的类型视为一个以类型为指标的类型族 $P:\type\to\type$，并将给定的对象看作 $P(\nat)$ 中的一个元素，然后沿着道路 $\id \nat\natp$ 进行输运即可。
这涉及的额外开销要小得多。

为简单起见，我们是在 \nat 和 \natp 两个类型的定义\emph{看起来相同}的情况下描述这一方法的。
然而，实践中更常见的情形是，两个定义并非字面上相同，但其中一个归纳原理蕴含另一个。
\index{type!unit}%
\index{natural numbers!encoded as list(unit)@encoded as $\lst\unit$}%
例如，考虑单元素类型上的列表类型 $\lst\unit$，它由以下两者生成：

\begin{itemize}
\item 一个元素 $\nil:\lst\unit$，以及
\item 一个函数 $\cons:\unit \times \lst\unit \to\lst\unit$。
\end{itemize}

这与 \nat 的定义并不相同，也不会产生相同的归纳原理。
$\lst\unit$ 的归纳原理是说，对于任意 $E:\lst\unit\to\type$ 以及递归式数据 $e_\nil:E(\nil)$ 和 $e_\cons : \prd{u:\unit}{\ell:\lst\unit} E(\ell) \to E(\cons(u,\ell))$，存在 $f:\prd{\ell:\lst\unit} E(\ell)$ 使得 $f(\nil)\jdeq e_\nil$ 且 $f(\cons(u,\ell))\jdeq e_\cons(u,\ell,f(\ell))$。
（我们将在\cref{sec:strictly-positive}看到如何推导一个归纳定义的归纳原理。）

现在假设我们定义 $0'' \defeq \nil: \lst\unit$，并用 $\suc''(\ell) \defeq \cons(\ttt,\ell)$ 定义 $\suc'':\lst\unit\to\lst\unit$。
那么，对于任意 $E:\lst\unit\to\type$ 以及 $e_0:E(0'')$ 和 $e_s:\prd{\ell:\lst\unit} E(\ell) \to E(\suc''(\ell))$，我们可以定义
\begin{align*}
  e_\nil &\defeq e_0\\
  e_\cons(\ttt,\ell,x) &\defeq e_s(\ell,x).
\end{align*}
（在定义 $e_\cons$ 时，我们利用了 \unit 的归纳原理来假设 $u$ 就是 $\ttt$。）
现在我们可以应用 $\lst\unit$ 的归纳原理，得到 $f:\prd{\ell:\lst\unit} E(\ell)$ 使得
\begin{gather*}
  f(0'') \jdeq f(\nil) \jdeq e_\nil \jdeq e_0\\
  f(\suc''(\ell)) \jdeq f(\cons(\ttt,\ell)) \jdeq e_\cons(\ttt,\ell,f(\ell)) \jdeq e_s(\ell,f(\ell)).
\end{gather*}
因此，$\lst\unit$ 满足与 $\nat$ 相同的归纳原理，从而（由与上述相同的论证）等于 $\nat$。

最后，这些结论并不局限于自然数：它们适用于任何归纳类型。
如果我们有一个归纳定义的类型 $W$，以及另一个类型 $W'$，满足与 $W$ 相同的归纳原理，那么可以推出 $\eqv{W}{W'}$，从而 $W=W'$。
我们分别利用 $W$ 和 $W'$ 导出的递归原理来构造映射 $W\to W'$ 和 $W'\to W$，然后利用各自的归纳原理证明两个复合都等于恒等映射。
例如，在\cref{cha:typetheory}中我们看到，余积 $A+B$ 也可以定义为$\sm{x:\bool} \rec{\bool}(\UU,A,B,x)$。
后者满足与前者相同的归纳原理；因此它们是典范等价的。

当然，这非常类似于范畴论中一个熟悉的事实：如果两个对象具有相同的\emph{泛性质}，那么它们等价。
在\cref{sec:initial-alg}中我们将看到，归纳类型实际上确实具有一个泛性质，因此这是那个一般原理的一个体现。
\index{universal property}

\section{\texorpdfstring{$\w$}{W}-类型}
\label{sec:w-types}

归纳类型的概念非常宽泛，这固然有利于它们的有用性和适用性，但也使得它们难以作为一个整体来研究。
幸运的是，它们都可以形式化地归约到几个特殊情形。
讨论这一归约过程超出了本书的范围——反正它对实践中使用类型论的数学家来说也无关紧要——但我们将花少许时间讨论一个我们尚未遇到的基本特殊情形。
这些就是 Martin-L{\"o}f 的\emph{$\w$-类型（$\w$-types）}，也称为\emph{良基树（well-founded trees）类型}。
\index{tree, well-founded}%
$\w$-类型是自然数、列表和二叉树等类型的推广，它足够宽泛，能涵盖\emph{任意}归纳类型的``递归''层面。

一个特定的 $\w$-类型通过给出两个参数 $A : \type$ 和 $B : A \to \type$ 来确定，在这种情况下，所得的 $\w$-类型记作 $\wtype{a:A} B(a)$。
类型 $A$ 表示 $\wtype{a :A} B(a)$ 的\emph{标签（labels）}类型，其作用类似于构造子（不过，我们将``构造子''一词保留给归纳定义中实际产生的函数）。
例如，当用 $\w$-类型定义自然数时，%
\index{natural numbers!encoded as a W-type@encoded as a $\w$-type}
类型 $A$ 将是包含两个元素 $\bfalse$ 和 $\btrue$ 的类型 $\bool$，因为恰好有两种方式得到一个自然数——它要么是零，要么是另一个自然数的后继。

类型族 $B : A \to \type$ 用于记录标签的\emph{元数（arity）}\index{arity}：一个标签 $a : A$ 将接受一族以 $B(a)$ 为指标的归纳参数。
因此，我们可以把 $B(a)$ 看作是标签 $a$ 的参数``个数''。
这些参数通过函数 $f : B(a) \to \wtype{a :A} B(a)$ 表示，其含义是，对于任何 $b : B(a)$，$f(b)$ 是该标签 $a$ 的``第 $b$ 个''参数。
因此，$\w$-类型 $\wtype{a :A} B(a)$ 可以被看作是良基树的类型，其中节点由 $A$ 中的元素标记，并且每个标记为 $a : A$ 的节点都有 $B(a)$ 条分支。

在自然数的例子中，标签 $\bfalse$ 的元数（arity）为 $0$，因为它构造的是常量零\index{zero}；标签 $\btrue$ 的元数为 $1$，因为它构造的是其参数的后继\index{successor}。
我们可以通过对 $\bool$ 做简单的消去来刻画这一点，定义一个到类型宇宙的函数 $\rec\bool(\bbU,\emptyt,\unit)$；对于 $\bfalse$，该函数返回空类型 $\emptyt$，对于 $\btrue$，则返回单位\index{type!unit}类型 $\unit$。
因此我们可以定义
\symlabel{natw}
\index{type!unit}%
\index{type!empty}%
\[ \natw \defeq \wtype{b:\bool} \rec\bool(\bbU,\emptyt,\unit,b) \]
其中上标 $\mathbf{w}$ 用于区分这个版本的自然数与之前使用的版本。
类似地，我们可以将 $A$ 上的列表\index{type!of lists}类型定义为一个具有 $\unit + A$ 个标签的 $\w$-类型：一个零元标签对应空列表，加上每个 $a : A$ 的一个一元标签，对应将 $a$ 附加到列表头部：
\[ \lst A \defeq \wtype{x: \unit + A} \rec{\unit + A}(\bbU, \; \emptyt, \; \lamu{a:A} \unit,\; x). \]
%
\index{W-type@$\w$-type}%
\indexsee{type!W-@$\w$-}{$\w$-type}%
一般而言，由 $A : \type$ 和 $B : A \to \type$ 所指定的 $\w$-类型 $\wtype{x:A} B(x)$，是由以下构造子生成的归纳类型：

\begin{itemize}
\item \label{defn:supp}
  $\supp : \prd{a:A} \Big(B(a) \to \wtype{x:A} B(x)\Big) \to \wtype{x:A} B(x)$.
\end{itemize}

%
构造子 $\supp$（上确界\index{supremum!constructor of a W-type@constructor of a $\w$-type}的缩写）接受一个标签 $a : A$ 以及一个函数 $f : B(a) \to \wtype{x:A} B(x)$（表示 $a$ 的参数），并构造出 $\wtype{x:A} B(x)$ 的一个新元素。
利用前面对自然数的 $\w$-类型编码，我们可以定义例如
\begin{equation*}
\zerow \defeq \supp(\bfalse, \; \lamu{x:\emptyt} \rec\emptyt(\natw,x)).
\end{equation*}
换句话说，我们使用标签 $\bfalse$ 来构造 $\zerow$。
而 $\rec\bool(\bbU,\emptyt,\unit, \bfalse)$ 会计算为 $\emptyt$，这符合预期，因为 $\bfalse$ 是一个零元标签。
因此，我们需要构造一个函数 $f : \emptyt \to \natw$，它表示提供给 $\bfalse$ 的（零个）参数。这当然是平凡的，如下所示，对 $\emptyt$ 做简单的消去即可。
类似地，我们可以定义 $1^{\mathbf{w}}$ 和一个后继函数 $\sucw$
\begin{align*}
1^{\mathbf{w}} &\defeq \supp(\btrue, \; \lamu{x:\unit} 0^{\mathbf{w}}) \\
\sucw&\defeq \lamu{n:\natw} \supp(\btrue, \; \lamu{x:\unit} n).
\end{align*}

\index{induction principle!for W-types@for $\w$-types}%
对于 $\w$-类型，我们有如下的归纳原理：

\begin{itemize}
\item 要证明关于 $\w$-类型 $\wtype{x:A} B(x)$ 中\emph{所有}元素的某个命题 $E : \big(\wtype{x:A} B(x)\big) \to \type$，只需对形如 $\supp(a,f)$ 的元素证明它，同时假设它对所有满足 $b : B(a)$ 的 $f(b)$ 成立。
换言之，只需给出一个证明
\begin{equation*}
e : \prd{a:A}{f : B(a) \to \wtype{x:A} B(x)}{g : \prd{b : B(a)} E(f(b))} E(\supp(a,f))
\end{equation*}
\end{itemize}

\index{variable}%
变量 $g$ 代表了我们的归纳假设，即 $a$ 的所有参数都满足 $E$。
为了表述这一点，我们对 $B(a)$ 类型的所有元素进行量化，因为每个 $b : B(a)$ 对应 $a$ 的一个参数 $f(b)$。

对于编码为 $\w$-类型的自然数，我们如何定义函数 $\dbl$？
我们希望使用 $\natw$ 的递归原理，其值域是 $\natw$ 本身。
因此需要构造一个合适的函数
\[e : \prd{a : \bool}{f : B(a) \to \natw}{g : B(a) \to \natw} \natw\]
它将代表 $\dbl$ 函数的递归式；为简便起见，我们将类型族 $\rec\bool(\bbU,\emptyt,\unit)$ 记作 $B$。

显然，$e$ 将是一个以 $a : \bool$ 作为其第一个参数的函数。
下一步是对 $a$ 进行分情形讨论，根据它是 $\bfalse$ 还是 $\btrue$ 来分别处理。
这提示了如下的形式
\[ e \defeq \lamu{a:\bool} \ind\bool(C,e_0,e_1,a) \]
其中
\[C \defeq \lamu{a:\bool} \prd{f : B(a) \to \natw}{g : B(a) \to \natw} \natw.\]
如果 $a$ 是 $\bfalse$，则类型 $B(a)$ 变为 $\emptyt$。
因此，给定 $f : \emptyt \to \natw$ 和 $g : \emptyt \to \natw$，我们希望构造一个 $\natw$ 的元素。
由于标签 $\bfalse$ 代表 $\emptyt$，它需要零个归纳参数，变量 $f$ 和 $g$ 无关紧要。
我们返回 $\zerow$\index{zero} 作为结果：
\[ e_0 \defeq \lamu{f:\emptyt \to \natw}{g:\emptyt \to \natw} \zerow. \]
类似地，如果 $a$ 是 $\btrue$，则类型 $B(a)$ 变为 $\unit$。
由于标签 $\btrue$ 代表后继\index{successor}运算，它需要一个归纳参数——即前驱\index{predecessor}——由变量 $f : \unit \to \natw$ 表示。
在前驱上的递归调用的值由变量 $g : \unit \to \natw$ 表示。
于是，取这个值（即 $g(\ttt)$）并应用后继函数两次，便得到想要的结果：
\begin{equation*}
e_1 \defeq \; \lamu{f:\unit \to \natw}{g:\unit \to \natw} \sucw(\sucw(g(\ttt))).
\end{equation*}
综上，我们得到
\[ \dbl \defeq \rec\natw(\natw, e) \]
其中 $e$ 如上所定义。

\symlabel{defn:recursor-wtype}
对于函数 $\rec{\wtype{x:A} B(x)}(E,e) : \prd{w : \wtype{x:A} B(x)} E(w)$，其对应的计算规则如下。
\index{computation rule!for W-types@for $\w$-types}%

\begin{itemize}
\item
  对于任意 $a : A$ 和 $f : B(a) \to \wtype{x:A} B(x)$，有
  \begin{equation*}
    \rec{\wtype{x:A} B(x)}(E,e,\supp(a,f)) \jdeq
    e(a,f,\big(\lamu{b:B(a)} \rec{\wtype{x:A} B(x)}(E,e,f(b))\big)).
  \end{equation*}
\end{itemize}

换言之，函数 $\rec{\wtype{x:A} B(x)}(E,e)$ 满足递归式 $e$。

根据上述计算规则，函数 $\dbl$ 的行为符合预期：
\begin{align*}
\dbl(\zerow) & \jdeq \rec\natw(\natw, e, \supp(\bfalse, \; \lamu{x:\emptyt} \rec\emptyt(\natw,x))) \\
& \jdeq e(\bfalse, \big(\lamu{x:\emptyt} \rec\emptyt(\natw,x)\big),
   \big(\lamu{x:\emptyt} \dbl(\rec\emptyt(\natw,x))\big)) \\
 & \jdeq e_0(\big(\lamu{x:\emptyt} \rec\emptyt(\natw,x)\big), \big(\lamu{x:\emptyt} \dbl(\rec\emptyt(\natw,x))\big))\\
 & \jdeq \zerow \\
 \intertext{and}
\dbl(1^{\mathbf{w}}) & \jdeq \rec\natw(\natw, e, \supp(\btrue, \; \lamu{x:\unit} \zerow)) \\
& \jdeq e(\btrue, \big(\lamu{x:\unit} \zerow\big), \big(\lamu{x:\unit} \dbl(\zerow)\big)) \\
 & \jdeq e_1(\big(\lamu{x:\unit} \zerow\big), \big(\lamu{x:\unit} \dbl(\zerow)\big)) \\
 & \jdeq \sucw(\sucw\big((\lamu{x:\unit} \dbl(\zerow))(\ttt)\big))\\
 & \jdeq \sucw(\sucw(\zerow))
\end{align*}
依此类推。

正如对于自然数一样，我们也能证明 $\w$-类型的唯一性定理：

\begin{thm}\label{thm:w-uniq}
  \index{uniqueness!principle, propositional!for functions on W-types@for functions on $\w$-types}%
设 $g,h : \prd{w:\wtype{x:A}B(x)} E(w)$ 为两个满足递归式
%
\begin{equation*}
  e : \prd{a,f} \Parens{\prd{b : B(a)} E(f(b))} \to  E(\supp(a,f)),
\end{equation*}
%
的函数，即在命题意义上满足该递归式，亦即
%
\begin{gather*}
 \prd{a,f} \id{g(\supp(a,f))} {e(a,f,\lamu{b:B(a)} g(f(b)))}, \\
 \prd{a,f} \id{h(\supp(a,f))}{e(a,f,\lamu{b:B(a)} h(f(b)))}.
\end{gather*}
那么 $g$ 与 $h$ 相等。
\end{thm}

\section{归纳类型是初始代数}
\label{sec:initial-alg}

\indexsee{initial!algebra characterization of inductive types}{homotopy-initial}%

如前所述，归纳类型也具有范畴论意义上的泛性质。
它们是\emph{同伦初始代数（homotopy-initial algebras）}：在由指定的构造子决定的``代数''范畴中，它们是在同伦相干的意义下的初始对象。
举一个简单的例子，考虑自然数。
这里合适的``代数''是指一个类型，它配备了与 $\nat$ 的构造子所赋予的相同结构。

\index{natural numbers!as homotopy-initial algebra}


\begin{defn}\label{defn:nalg}
  一个\define{$\nat$-代数（$\nat$-algebra）}
  \indexdef{N-algebra@$\nat$-algebra}%
  \indexdef{algebra!N-@$\nat$-}%
  是一个类型 $C$ 连同两个元素 $c_0 : C$ 与 $c_s : C \to C$。这类代数的类型为
\begin{equation*}
\nalg \defeq \sm {C : \type} C \times (C \to C).
\end{equation*}
\end{defn}

\begin{defn}\label{defn:nhom}
  在两个 $\nat$-代数 $(C,c_0,c_s)$ 与 $(D,d_0,d_s)$ 之间的一个\define{$\nat$-同态（$\nat$-homomorphism）}
  \indexdef{N-homomorphism@$\nat$-homomorphism}%
  \indexdef{homomorphism!N-@$\nat$-}%
  是一个函数 $h : C \to D$，满足 $h(c_0) = d_0$ 且对所有 $c : C$ 有 $h(c_s(c)) = d_s(h(c))$。这类同态的类型为
\begin{narrowmultline*}
\nhom((C,c_0,c_s),(D,d_0,d_s)) \defeq \narrowbreak
 \dsm {h : C \to D} (\id{h(c_0)}{d_0}) \times \tprd{c:C} (\id{h(c_s(c))}{d_s(h(c))}).
\end{narrowmultline*}
\end{defn}

这样我们就得到了一个由 $\nat$-代数与 $\nat$-同态构成的范畴，而我们的断言是：$\nat$ 是这个范畴的初始对象。
范畴论学者会立刻认出，这就是范畴中\emph{自然数对象（natural numbers object）}的定义。

当然，由于我们的类型表现得像 $\infty$-群胚\index{.infinity-groupoid@$\infty$-groupoid}，我们实际上有一个 $\nat$-代数的 $(\infty,1)$-范畴\index{.infinity1-category@$(\infty,1)$-category}，因此我们应该要求 $\nat$ 在适当的 $(\infty,1)$-范畴意义下是初始的。
幸运的是，我们无需定义 $(\infty,1)$-范畴也能表述这一点。

\begin{defn}
  \index{universal!property!of natural numbers}%
  一个 $\nat$-代数 $I$ 被称为\define{同伦初始的（homotopy-initial）}，
  \indexdef{homotopy-initial!N-algebra@$\nat$-algebra}%
  \indexdef{N-algebra@$\nat$-algebra!homotopy-initial (h-initial)}%
  或简称为\define{h-初始的（h-initial）}，
  \indexsee{h-initial}{homotopy-initial}%
  如果对于任何其他 $\nat$-代数 $C$，从 $I$ 到 $C$ 的 $\nat$-同态类型是可缩的。也就是说，
\begin{equation*}
\ishinitn(I) \defeq \prd{C : \nalg} \iscontr(\nhom(I,C)).
\end{equation*}
\end{defn}

当 h-初始代数存在时，它们是唯一的——不仅是范畴论中常见的``在同构意义下''唯一，更由泛等公理在``相等''的意义下唯一。

\begin{thm}
  任何两个 h-初始的 $\nat$-代数都是相等的。
  因此，h-初始 $\nat$-代数的类型是一个命题。
\end{thm}

\begin{proof}
  设 $I$ 和 $J$ 是 h-初始的 $\nat$-代数。
  那么 $\nhom(I,J)$ 是可缩的，因而其中存在某个 $\nat$-同态 $f:I\to J$；同理我们也有一个 $\nat$-同态 $g:J\to I$。
  现在复合 $g\circ f$ 是一个从 $I$ 到 $I$ 的 $\nat$-同态，$\idfunc[I]$ 也是；而 $\nhom(I,I)$ 是可缩的，所以 $g\circ f = \idfunc[I]$。
  类似地，$f\circ g = \idfunc[J]$。
  因此 $\eqv IJ$，从而 $I=J$。由于``可缩''是一个命题，并且依值函数类型保持命题，所以``是 h-初始的''本身也是一个命题。因此，关于 $I$（或 $J$）为 h-初始的任何两个证明都必然相等，这就完成了证明。
\end{proof}

我们现在有以下定理。

\begin{thm}\label{thm:nat-hinitial}
$\nat$-代数 $(\nat, \emptyt, \suc)$ 是同伦初始的。
\end{thm}

\begin{proof}[证明概要]
  固定任意一个 $\nat$-代数 $(C,c_0,c_s)$。
  $\nat$ 的递归原理给出一个函数 $f:\nat\to C$，定义为
  \begin{align*}
    f(0) &\defeq c_0\\
    f(\suc(n)) &\defeq c_s(f(n)).
  \end{align*}
  这两个等式使得 $f$ 成为一个 $\nat$-同态，我们可以将其取为 $\nhom(\nat,C)$ 的收缩中心。
  唯一性定理（\cref{thm:nat-uniq}）则意味着任何其他 $\nat$-同态都等于 $f$。
\end{proof}

为了将其置于更一般的上下文中，考虑\emph{自函子的代数（algebra for an endofunctor）}这一概念会有所帮助。\index{algebra!for an endofunctor}
注意，将一个类型 $C$ 变成 $\nat$-代数，等同于给出一个函数 $c:C+\unit\to C$；而一个函数 $f:C\to D$ 是一个 $\nat$-同态，当且仅当 $f \circ c \htpy d \circ (f+\unit)$。
用范畴论的语言来说，这意味着 $\nat$-代数就是类型范畴上的自函子 $F(X)\defeq X+1$ 的代数。

\indexsee{functor!polynomial}{endofunctor, polynomial}%
\indexsee{polynomial functor}{endofunctor, polynomial}%
\indexdef{endofunctor!polynomial}%
\index{W-type@$\w$-type!as homotopy-initial algebra}
考虑一个更一般的情形，对于给定的 $A : \type$ 和 $B : A \to \type$，考察与之关联的 $\w$ 类型。
在这种情况下，我们有一个关联的\define{多项式函子}：
\begin{equation}
\label{eq:polyfunc}
P(X) = \sm{x : A} (B(x) \rightarrow X).
\end{equation}
实际上，这个指派只是在同伦意义下具有函子性，但这不影响接下来的讨论。
根据定义，一个\define{$P$-代数（$P$-algebra）}
\indexdef{algebra!for a polynomial functor}%
\indexdef{algebra!W-@$\w$-}%
就是一个类型 $C$ 配上一个函数 $s_C :  PC \rightarrow C$。
由 $\Sigma$ 类型的泛性质，这等价于给出一个函数 $\prd{a:A} (B(a) \to C) \to C$。
我们也称这样的对象为关于 $A$ 和 $B$ 的\define{$\w$-代数（$\w$-algebra）}，
\indexdef{W-algebra@$\w$-algebra}%
并记
\symlabel{walg}
\begin{equation*}
\walg(A,B) \defeq \sm {C : \type} \prd{a:A} (B(a) \to C) \to C.
\end{equation*}

类似地，对于 $P$-代数 $(C,s_C)$ 和 $(D,s_D)$，它们之间的一个\define{同态}
\indexdef{homomorphism!of algebras for a functor}%
$(f, s_f) : (C, s_C) \rightarrow (D, s_D)$ 由一个函数 $f : C \rightarrow D$ 和映射 $PC \rightarrow D$ 之间的一个同伦构成：
\[
s_f :  f \circ s_C \, = s_{D} \circ Pf,
\]
其中 $Pf : PC\rightarrow PD$ 是 $P$ 对 $f: C \rightarrow D$ 施加作用的结果（该作用易于定义）。这样的代数同态可以示意性地表示为：
\[
\xymatrix{
 PC \ar[d]_{s_C} \ar[r]^{Pf}  \ar@{}[dr]|{s_f} &  PD \ar[d]^{s_D}\\
C \ar[r]_{f}   & D }
\]
就元素而言，$f$ 是一个 $P$-同态（或称为\define{$\w$-同态（$\w$-homomorphism）}）当且仅当
\indexdef{W-homomorphism@$\w$-homomorphism}%
\indexdef{homomorphism!W-@$\w$-}%
\[f(s_C(a,h)) = s_D(a,f \circ h).\]
我们有 $\w$-同态的类型：
\symlabel{whom}
\begin{equation*}
  \whom_{A,B}((C, s_C),(D,s_D)) \defeq \sm{f : C \to D} \prd{a:A}{h:B(a)\to C} \id{f(s_C(a,h))}{s_D(a, f\circ h)}
\end{equation*}

\index{universal!property!of $\w$-type}%
最后，一个 $P$-代数 $(C, s_C)$ 称为\define{同伦初始的}
\indexdef{homotopy-initial!algebra for a functor}%
\indexdef{homotopy-initial!W-algebra@$\w$-algebra}%
，如果对于每一个 $P$-代数 $(D, s_D)$，所有代数同态的类型 $(C, s_C) \rightarrow (D, s_D)$ 是可缩的。
也就是说，
\begin{equation*}
\ishinitw(A,B,I) \defeq \prd{C : \walg(A,B)} \iscontr(\whom_{A,B}(I,C)).
\end{equation*}
%
现在，与 \cref{thm:nat-hinitial} 类似的定理是：

\begin{thm}\label{thm:w-hinit}
对于任意类型 $A : \type$ 和类型族 $B : A \to \type$，$\w$-代数 $(\wtype{x:A}B(x), \supp)$ 是同伦初始的。
\end{thm}

\begin{proof}[证明概要]
假设我们有 $A : \type$ 和 $B : A \to \type$，
并考虑关联的多项式函子 $P(X)\defeq\sm{x:A}(B(x)\to X)$。
令 $W \defeq \wtype{x:A}B(x)$。那么利用
\cref{sec:w-types} 中的 $\w$ 引入规则，我们有一个结构映射 $s_W\defeq\supp: PW \rightarrow W$。
我们希望证明代数 $(W, s_W)$ 是
同伦初始的。因此，考虑另一个代数 $(C,s_C)$，并证明从 $(W, s_W)$ 到 $(C, s_C)$ 的 $\w$-同态的类型 $T\defeq \whom_{A,B}((W, s_W),(C,s_C)) $
是可缩的。为此，
观察到 $\w$-消去规则和 $\w$-计算
规则允许我们定义一个 $\w$-同态 $(f, s_f) : (W, s_W) \rightarrow (C, s_C)$，
从而证明 $T$ 是有实例的。此外还需要证明，对于每一个 $\w$-同态 $(g, s_g) : (W, s_W) \rightarrow (C, s_C)$，存在一个恒等证明
\begin{equation}
\label{equ:prequired}
p :  (f,s_f) = (g,s_g).
\end{equation}
这里用到如下一般性事实：形如 $(f,s_f) = (g,s_g) $ 的类型
等价于我们所谓的从 $f$ 到 $g$ 的 \define{代数 $2$-胞（algebra $2$-cells）}
\indexdef{algebra!2-cell}%
的类型，其典范
元素形如 $(e, s_e)$，其中 $e : f=g$，而 $s_e$ 是一个更高阶的恒等证明，连接如下粘合图表所表示的恒等证明：
\[
\xymatrix{
PW \ar@/^1pc/[r]^{Pg}   \ar[d]_{s_W}   \ar@/_1pc/[r]_{Pf} \ar@{}[r]|{Pe}
& PC \ar[d]^{s_C}  \\
W  \ar@/_1pc/[r]_f  \ar@{}[r]^{s_f} & C } \qquad
\xymatrix{
PW \ar@/^1pc/[r]^{Pg}   \ar[d]_{s_W} \ar@{}[r]_(.52){s_g}  & PC \ar[d]^{s_C}  \\
W \ar@/^1pc/[r]^g  \ar@/_1pc/[r]_f  \ar@{}[r]|{e} & C }
\]
鉴于这一事实，要证明存在如~\eqref{equ:prequired} 所示的元素，只需
证明存在一个从 $f$ 到 $g$ 的代数 2-胞 $(e,s_e)$ 即可。
恒等证明 $e : f=g$ 现在可以通过函数外延性和
$\w$-消去规则来构造，以保证所需的恒等
证明 $s_e$ 的存在性。
\end{proof}

%%%%%%%%%%%%%%%%%%%%%%%%%%%%%%%%%%%%%%%%%%%%%%%%%%%%%%%%%%

\section{同伦归纳类型}
\label{sec:htpy-inductive}

在 \cref{sec:w-types} 中，我们展示了如何将自然数编码为 $\w$ 类型，即
\begin{align*}
\natw & \defeq \wtype{b:\bool} \rec\bool(\bbU,\emptyt,\unit,b), \\
\zerow & \defeq \supp(\bfalse, (\lamu{x:\emptyt} \rec\emptyt(\natw,x))), \\
\sucw & \defeq \lamu{n:\natw} \supp(\btrue, (\lamu{x:\unit} n)).
\end{align*}
我们还展示了如何利用递归原理在 $\natw$ 上定义 $\dbl$ 函数。
然而，当涉及到归纳原理时，这种编码方式就不再令人满意了：给定 $E : \natw \to \type$ 以及递归式 $e_z : E(\zerow)$ 和 $e_s : \prd{n : \natw}  E(n) \to E(\sucw(n))$，我们只能构造一个依值函数 $r(E,e_z,e_s) : \prd{n : \natw} E(n)$，它\emph{仅在命题意义上}满足给定的递归式，即只相差一条道路。
这意味着，自然数的计算规则（它给出的是判断相等）无法从 $\w$ 类型的规则中以任何明显的方式推导出来。

\index{type!homotopy-inductive}%
\index{homotopy-inductive type}%
如果我们不考虑传统的归纳类型，而是考虑\emph{同伦归纳类型（homotopy-inductive types）}，那么这个问题就不复存在了。在同伦归纳类型中，所有计算规则都表述为道路层面的相等，也就是说，符号 $\jdeq$ 被替换为 $=$。例如，同伦版本 $\w$ 类型 $\mathsf{W^h}$ 的计算规则变为：
\index{computation rule!propositional}%


\begin{itemize}
\item 对于任意 $a : A$ 与 $f : B(a) \to \wtypeh{x:A} B(x)$，我们有
\begin{equation*}
  \rec{\wtypeh{x:A} B(x)}(E,e,\supp(a,f)) = e\Big(a,f,\big(\lamu{b:B(a)} \rec{\wtypeh{x:A} B(x)}(E,f(b))\big)\Big)
\end{equation*}
\end{itemize}

同伦归纳类型在计算性质上有一个明显的劣势——任何使用归纳原理构造的函数，其行为现在都只能命题地表征。
但仍有诸多其他考量促使我们去研究同伦归纳类型。
例如，尽管我们在 \cref{sec:initial-alg} 中说明了归纳类型是同伦初始代数，但并非每个同伦初始代数都是归纳类型（即满足相应的归纳原理）——而每个同伦初始代数\emph{都}是一个同伦归纳类型。
类似地，当涉及的类型中有一个（或两个）只是同伦归纳类型时，我们可能仍想应用来自 \cref{sec:appetizer-univalence} 的唯一性论证——比方说，用于证明用 $\w$-类型编码的 $\nat$ 等价于通常的 $\nat$。

此外，同伦归纳类型这一概念现在是类型论\emph{内部}的了。
例如，这意味着我们可以形成由所有自然数对象构成的类型，并对其作出断言。
对于 $\w$-类型的情形，我们可以将同伦 $\w$-类型 $\wtype{x:A} B(x)$ 刻画为：配备了上确界函数以及满足恰当（命题的）计算规则的归纳原理的任何类型：
\begin{multline*}
\w_d(A,B) \defeq \sm{W : \type}
                 \sm{\supp : \prd {a} (B(a) \to W) \to W}
                 \prd{E : W \to \type} \\
                 \prd{e : \prd{a,f} (\prd{b : B(a)} E(f(b))) \to E(\supp(a,f))}
                 \sm{\ind{} : \prd{w : W} E(w)}
                 \prd{a,f} \\
                 \ind{}(\supp(a,f)) = e(a,\lamu{b:B(a)} \ind{}(f(b))).
\end{multline*}
在 \cref{cha:hits} 中，我们将看到命题计算规则值得考虑的其他一些原因。

本节中，我们将陈述关于同伦归纳类型的一些基本事实。
我们省略大部分证明，因为它们技术性较强。

\begin{thm}
  对于任意 $A : \type$ 与 $B : A \to \type$，类型 $\w_d(A,B)$ 是一个命题。
\end{thm}

事实上，存在一种使用递归原理外加某些\emph{唯一性}与\emph{相干性}定律的等价刻画。首先给出递归原理：
%


\begin{itemize}
\item 当构造从 $\w$-类型 $\wtypeh{x:A} B(x)$ 到类型 $C$ 的函数时，只要对 $\supp(a,f)$ 给出其值就足够了——这里假设我们已经得到了所有 $b : B(a)$ 对应的 $f(b)$ 的值。
换言之，只需构造一个函数
\begin{equation*}
  c : \prd{a:A} (B(a) \to C) \to C.
\end{equation*}
\end{itemize}


\index{computation rule!propositional}%
对于 $\rec{\wtypeh{x:A} B(x)}(C,c) : (\wtype{x:A} B(x)) \to C$ 的相应计算规则如下：


\begin{itemize}
\item 对于任意 $a : A$ 与 $f : B(a) \to \wtypeh{x:A} B(x)$，我们有一个见证 $\beta(C,c,a,f)$，表明下述等式成立：
\begin{equation*}
  \rec{\wtypeh{x:A} B(x)}(C,c,\supp(a,f)) =
  c(a,\lamu{b:B(a)} \rec{\wtypeh{x:A} B(x)}(C,c,f(b))).
\end{equation*}
\end{itemize}

此外，我们断言下述唯一性原理，它表明任何两个由相同递归式定义的函数都是相等的：
\index{uniqueness!principle, propositional!for homotopy W-types@for homotopy $\w$-types}%


\begin{itemize}
\item 设给定 $C : \type$ 与 $c : \prd{a:A} (B(a) \to C) \to C$。令 $g,h : (\wtypeh{x:A} B(x)) \to C$ 为两个函数，它们满足命题相等意义下的递归式 $c$，即我们有
\begin{align*}
  \beta_g &: \prd{a,f} \id{g(\supp(a,f))}{c(a,\lamu{b: B(a)} g(f(b)))}, \\
  \beta_h &: \prd{a,f} \id{h(\supp(a,f))}{c(a,\lamu{b: B(a)} h(f(b)))}.
\end{align*}
那么 $g$ 与 $h$ 相等，即存在一个 $\alpha(C,c,f,g,\beta_g,\beta_h)$，其类型为 $g = h$。
\end{itemize}

\index{coherence}%
回想一下，当我们拥有的是归纳原理而不仅仅是递归原理时，这个命题唯一性原理（propositional uniqueness principle）是可以导出的（\cref{thm:w-uniq}）。
但如果只有递归原理，唯一性原理就不再可导出了——事实上，这条陈述甚至都不成立（练习）。因此，我们将其假定为一条公理。
我们还假定了以下相干性（coherence）\index{coherence}定律，它告诉我们唯一性的证明在典范元上的行为方式：


\begin{itemize}
\item
对于任意的 $a : A$ 和 $f : B(a) \to C$，下述图表在命题意义下交换：
\[\xymatrix{
  g(\supp(a,f)) \ar_{\alpha(\supp(a,f))}[d] \ar^-{\beta_g}[r] & c(a,\lamu{b:B(a)} g(f(b)))
  \ar^{c(a,\blank)(\funext (\lam{b} \alpha(f(b))))}[d] \\
  h(\supp(a,f)) \ar_-{\beta_h}[r] & c(a,\lamu{b: B(a)} h(f(b))) \\
}\]
其中 $\alpha$ 是道路 $\alpha(C,c,f,g,\beta_g,\beta_h) : g = h$ 的缩写。
\end{itemize}

将这些数据综合起来，便得到了 $\wtype{x:A} B(x)$ 的另一种刻画：它是一个具有上确界函数的类型，满足简单的消去、计算、唯一性与相干性规则：
\begin{multline*}
\w_s(A,B) \defeq \sm{W : \type}
                       \sm{\supp : \prd {a} (B(a) \to W) \to W}
                       \prd{C : \type}
                       \prd{c : \prd{a} (B(a) \to C) \to C}\\
                       \sm{\rec{} : W \to C}
                       \sm{\beta : \prd{a,f} \rec{}(\supp(a,f)) = c(a,\lamu{b: B(a)} \rec{}(f(b)))} \narrowbreak
                       \prd{g : W \to C}
                       \prd{h : W \to C}
                       \prd{\beta_g : \prd{a,f} g(\supp(a,f)) = c(a,\lamu{b: B(a)} g(f(b)))} \\
                       \prd{\beta_h : \prd{a,f} h(\supp(a,f)) = c(a,\lamu{b: B(a)} h(f(b)))}
                       \sm{\alpha : \prd {w : W} g(w) = h(w)}
                       \prd{a,f} \\
                       \alpha(\supp(a,f)) \ct \beta_h = \beta_g \ct c(a,-)(\funext \; \lam{b} \alpha(f(b)))
\end{multline*}

\begin{thm}
对于任意 $A : \type$ 和 $B : A \to \type$，类型 $\w_s (A,B)$ 是一个命题。
\end{thm}

最后，我们还有第三种非常简洁的刻画：将 $\wtype{x:A} B(x)$ 视为一个同伦初始的 $\w$-代数：
\begin{equation*}
\w_h(A,B) \defeq \sm{I : \walg(A,B)} \ishinitw(A,B,I).
\end{equation*}

\begin{thm}
对于任意 $A : \type$ 和 $B : A \to \type$，类型 $\w_h (A,B)$ 是一个命题。
\end{thm}

事实上，$\w$-类型的这三种刻画是等价的：


\begin{lem}\label{lem:homotopy-induction-times-3}
对于任意 $A : \type$ 和 $B : A \to \type$，我们有
\[ \w_d(A,B) \eqvsym \w_s(A,B) \eqvsym \w_h(A,B) \]
\end{lem}

更进一步，我们有如下定理，它改进了 \cref{thm:w-hinit}：

\begin{thm}
满足 $\w$-类型的形成、引入、消去与命题计算规则的类型，正是同伦初始的 $\w$-代数。
\end{thm}

%%%%%

\begin{proof}[证明概要]
%%%%%
回顾 \cref{thm:w-hinit} 的证明可知，要确立 $\wtype{x:A}B(x)$ 的同伦初始性（h-initial），仅仅用到其\emph{命题}计算规则就已足够。
为证其逆，假设与 $A : \type$ 和 $B : A \to \UU$ 关联的多项式函子（polynomial functor）有一个同伦初始代数 $(W,s_W)$；我们来证明 $W$ 满足 $\w$-类型的命题规则。
$\w$-引入规则很简单：对于 $a : A$ 和 $t : B(a) \rightarrow W$，我们定义 $\supp(a,t) : W$ 为将结构映射 $s_W : PW \rightarrow W$ 应用于 $(a,t) : PW$ 的结果。
对于 $\w$-消去规则，我们假设其前提成立，特别是有 $C' : W \to \type$。
利用其他前提，可以证明类型 $C \defeq \sm{ w : W} C'(w)$ 可以配备一个结构映射 $s_C : PC \rightarrow C$。根据 $W$ 的同伦初始性，我们得到一个代数同态 $(f, s_f) : (W, s_W) \rightarrow (C, s_C)$。此外，第一个投影 $\proj1 : C \rightarrow W$ 也能配备同态结构，因此我们得到如下形式的图表：
\[
\xymatrix{
PW \ar[r]^{Pf} \ar[d]_{s_W}  & PC \ar[d]^{s_C}  \ar[r]^{P \proj1}  & PW  \ar[d]^{s_W}  \\
W \ar[r]_f  & C \ar[r]_{\proj1}  & W.}
\]
但恒等函数 $1_W : W \rightarrow W$ 具有典范的代数同态结构，所以根据从 $(W,s_W)$ 到自身的同态类型的可缩性，必须存在一个恒等证明，表明 $(f,s_f)$ 与 $(\proj1, s_{\proj1})$ 的复合等于 $(1_W, s_{1_W})$。这特别意味着存在一个恒等证明 $p :  \proj1 \circ f = 1_W$。

由于 $(\proj2 \circ f) w : C( (\proj1 \circ f) w)$，我们可以定义
\[
\rec{}(w,c) \defeq
p_{\, * \,}( ( \proj2 \circ  f)   w )   : C(w)
\]
其中输运 $p_{\, * \,}$ 是关于族
\[
\lamu{u}C\circ u : (W\to W)\to W\to \UU.
\] 的。
命题 $\w$-计算规则的验证是一系列计算，涉及形如 $p_{\, * \,}$ 的操作的自然性性质。
\end{proof}


%%%%%

\index{natural numbers!encoded as a W-type@encoded as a $\w$-type}%
最后，如我们所愿，可以将同伦自然数编码为同伦 $\w$-类型：

\begin{thm}
带有命题计算规则的自然数规则，可以从带有命题计算规则的 $\w$-类型规则中推导出来。
\end{thm}

%%%%%%%%%%%%%%%%%%%%%%%%%%%%%%%%%%%%%%%%%%%%%%%%%%%%%%%%%%

\section{归纳定义的一般语法}
\label{sec:strictly-positive}

\index{type!inductive|(}%
\indexsee{inductive!type}{type, inductive}%

到目前为止，我们只讨论了特定的归纳类型：$\emptyt$、$\unit$、$\bool$、$\nat$、余积、积、$\Sigma$-类型、$\w$-类型等等。
然而，类型论的一个重要方面是能够定义\emph{新的}归纳类型，而不仅仅局限于某些特定的固定列表。
但要做到这一点，我们需要知道什么样的“归纳定义”是有效或合理的。

要知道并非所有“看起来像归纳定义”的东西都有意义，考虑下面这个类型 $C$ 的“构造子”：

\begin{itemize}
\item $g:(C\to \nat) \to C$.
\end{itemize}

此类类型 $C$ 的递归原理大概会说：给定一个类型 $P$，要构造一个函数 $f:C\to P$，只需考虑输入 $c:C$ 是形如 $g(\alpha)$ 的情况，其中 $\alpha:C\to\nat$。此外，我们期望能够以某种方式使用“递归数据”，即 $f$ 应用于 $\alpha$ 的结果。然而，“将 $f$ 应用于 $\alpha$”究竟是什么意思并不清楚，因为两者都是以 $C$ 为定义域的函数。

我们可以为 $C$ 写下一条“递归原理”，即（无根据地）假设存在某种方法能将 $f$ 应用于 $\alpha$ 并获得一个函数 $P\to\nat$。那么，递归规则的输入会要求一个类型 $P$ 连同函数
\begin{equation}
  h:(C\to\nat) \to (P\to\nat) \to P\label{eq:fake-recursor}
\end{equation}
其中 $h$ 的两个参数是 $\alpha$ 和“将 $f$ 应用于 $\alpha$ 的结果”。但是，所得的函数 $f:C\to P$ 的计算规则会是什么样呢？观察其他计算规则，我们会期望类似于“$f(g(\alpha)) \jdeq h(\alpha,f(\alpha))$”这样的东西，但如前所述，“$f(\alpha)$”没有意义。$C$ 的归纳原理就更成问题了；甚至连如何写下其前提都不清楚。

另一方面，我们可以为 $C$ 写下另一条不同的“递归原理”，即忽略 $\alpha$ 定义域中 $C$ 的“递归”存在，仅将其视为一簇自然数的指标类型。在这种情况下，输入会要求一个类型 $P$ 连同函数
\begin{equation*}
  h:(C\to \nat) \to P,
\end{equation*}
因此，递归原理的类型就是 $\rec{C}:\prd{P:\UU} ((C\to \nat) \to P) \to C\to P$，归纳原理也类似。现在可以写出一条计算规则，即 $\rec{C}(P,h,g(\alpha))\jdeq h(\alpha)$。然而，事实证明，存在一个具有此种递归子和计算规则的类型 $C$ 会导致矛盾。参见 \cref{ex:loop,ex:loop2,ex:inductive-lawvere,ex:ilunit} 了解其证明及其他变体。

上面的例子提示了归纳定义的一条限制：所有构造子的定义域都必须是所定义类型的\emph{协变函子}\index{functor!covariant}\index{covariant functor}，这样才能对它们“应用 $f$”，以得到“递归调用”的结果。
换句话说，如果把要定义的类型的所有出现都替换成一个变量
\index{variable!type}%
$X:\type$，那么每个构造子的定义域
\index{domain!of a constructor}%
必须是一个表达式，并且能够成为 $X$ 的一个协变函子。
我们目前考虑过的所有例子都满足这条限制。
例如，对于构造子 $\inl:A\to A+B$，对应的函子是取值于 $A$ 的常值函子（即 $X\mapsto A$）；
而对于构造子 $\suc:\nat\to\nat$，对应的函子是恒等函子（$X\mapsto X$）。

然而，这条必要条件并不是充分的。
协变性确实能防止归纳类型出现在单层函数类型的左侧，就像前面那个“构造子” $g$ 的参数 $C\to\nat$ 那样——因为那会产生一个逆变函子\index{functor!contravariant}\index{contravariant functor}，而不是协变函子。
但由于两个逆变函子的复合是协变的，因此像 $((X\to \nat)\to \nat)$ 这样的\emph{双重}函数类型又成为协变的。
这使我们能够重现 Cantor 风格的悖论\index{paradox}。

例如，考虑如下的“归纳类型” $D$，它具有如下构造子：

\begin{itemize}
\item $k:((D\to\prop)\to\prop)\to D$。
\end{itemize}

假设这样的类型存在，我们定义函数
\begin{align*}
  r&:D\to (D\to\prop)\to\prop,\\
  f&:(D\to\prop) \to D,\\
  p&:(D\to \prop) \to (D\to\prop)\to \prop,\\
  \intertext{by}
  r(k(\theta)) &\defeq \theta,\\
  f(\delta) &\defeq k(\lam{x} (x=\delta)),\\
  p(\delta) &\defeq \lam{x} \delta(f(x)).
\end{align*}
其中 $r$ 是利用 $D$ 的递归原理定义的，而 $f$ 和 $p$ 是显式定义的。
那么对任意 $\delta:D\to\prop$，有 $r(f(\delta)) = \lam{x}(x=\delta)$。

因此特别地，如果 $f(\delta)=f(\delta')$，那么我们有一条道路 $s:(\lam{x}(x=\delta)) = (\lam{x}(x=\delta'))$。
于是，$\happly(s,\delta) : (\delta=\delta) = (\delta=\delta')$，所以特别地有 $\delta=\delta'$ 成立。
从而，$f$ 是“单射”（尽管\emph{先验地} $D$ 可能不是一个集合）。
这已然颇为可疑——我们有了从 $D$ 的“幂集”\index{power set}到 $D$ 的“单射”——再稍加处理，便可将其化为矛盾。

假设给定 $\theta:(D\to\prop)\to\prop$，并定义 $\delta:D\to\prop$ 为
\begin{equation}
  \delta(d) \defeq \exis{\gamma:D\to\prop} (f(\gamma) = d) \times \theta(\gamma).\label{eq:Pinj}
\end{equation}
我们断言 $p(\delta)=\theta$。
根据函数外延性，只需证明对任意 $\gamma:D\to\prop$ 都有 $p(\delta)(\gamma) =_{\prop} \theta(\gamma)$。
而根据泛等性，为此只需证明两者互相蕴含即可。
现在由 $p$ 的定义，我们有
\begin{align*}
  p(\delta)(\gamma) &\jdeq \delta(f(\gamma))\\
  &\jdeq \exis{\gamma':D\to\prop} (f(\gamma') = f(\gamma)) \times \theta(\gamma').
\end{align*}
若 $\theta(\gamma)$ 成立，则此式显然成立，因为我们可以取 $\gamma'\defeq \gamma$。
另一方面，如果我们有 $\gamma'$ 满足 $f(\gamma') = f(\gamma)$ 且 $\theta(\gamma')$，那么由 $f$ 是单射可得 $\gamma'=\gamma$，因此也有 $\theta(\gamma)$。

这就完成了 $p(\delta)=\theta$ 的证明。
于是，每个元素 $\theta:(D\to\prop)\to\prop$ 都是某个元素 $\delta:D\to\prop$ 在 $p$ 下的像。
然而，如果我们用经典的对角线方法定义 $\theta$：
\[ \theta(\gamma) \defeq \neg p(\gamma)(\gamma) \quad\text{for all $\gamma:D\to\prop$} \]
那么由 $\theta = p(\delta)$ 可以推出 $p(\delta)(\delta) = \neg p(\delta)(\delta)$。
这是一个矛盾：没有一个命题可以等价于它的否定。
（假设 $P\Leftrightarrow \neg P$，如果 $P$，那么 $\neg P$，于是得到 $\emptyt$；因此 $\neg P$，但又推出 $P$，于是得到 $\emptyt$。）

\begin{rmk}
  这里需要处理宇宙大小的问题。
  一般来说，一个归纳类型必须存在于一个已经包含了其定义中所有类型的宇宙中。
  因此，如果在 $D$ 的定义中，模糊的记号 \prop 指的是 $\prop_{\UU}$，那么我们并没有 $D:\UU$，而只有对某个更大的满足 $\UU:\UU'$ 的宇宙 $\UU'$，有 $D:\UU'$。
  \index{mathematics!predicative}%
  \indexsee{impredicativity}{mathematics, predicative}%
  \indexsee{predicative mathematics}{mathematics, predicative}%
  因此，在直谓理论中，~\eqref{eq:Pinj} 的右侧属于 $\prop_{\UU'}$ 而非 $\prop_{\UU}$。
  所以这个矛盾的推导确实需要用到\cref{subsec:prop-subsets}中提到的命题大小重整公理。
  \index{propositional!resizing}%
\end{rmk}

\index{consistency}%
这个反例表明，我们应该禁止一个归纳类型出现在其构造子的定义域中的箭头左侧，即使这种出现嵌套在其他箭头中以至于最终是协变的。
（类似地，我们也禁止它出现在依值函数类型的定义域中。）
这种限制被称为\define{严格正性（strict positivity）}
\indexdef{strict!positivity}%
\indexsee{positivity, strict}{strict positivity}%
（普通的“正性”本质上是协变性），且事实证明它是充分的。

\index{constructor}%
总而言之，类型 $W$ 的一个有效归纳定义由一列\emph{构造子（constructor）}构成。
每个构造子具有一个函数类型：它取若干（可能为零）个输入（这些输入可能相互依赖），并返回 $W$ 的一个元素。
最后，我们允许 $W$ 自身出现在其构造子的输入类型中，但仅限于\emph{严格正的（strictly positive）}出现。
这实质上意味着，构造子的每个参数要么是一个不涉及 $W$ 的类型，要么是某种迭代的函数类型，且其最终的上域为 $W$。
例如，下面就是一个合法的构造子类型：
\begin{equation}
  c:(A\to W) \to (B\to C \to W) \to D \to W \to W.\label{eq:example-constructor}
\end{equation}
所有这些函数类型同样可以是依值函数（$\Pi$-类型）。%
\footnote{用 \cref{sec:initial-alg} 的语言来说，严格正性的条件确保了相关的自函子是多项式函子。\indexfoot{endofunctor!polynomial}\indexfoot{algebra!initial}\indexsee{algebra!initial}{homotopy-initial} 范畴论中众所周知，并非\emph{所有}自函子都具有初始代数；限制于多项式函子确保了相容性。 对这一条件可以有多种放宽方案，但在本书中，我们将限定于此处定义的严格正性。}

注意，我们要求归纳定义由\emph{有限}个构造子给出。
这仅仅是因为我们必须将其写在纸面上。
如果我们想要一个表现得仿佛拥有无限多个构造子的归纳类型，只需用某个无限类型来参数化一个构造子即可。
例如，像 $\nat \to W \to W$ 这样的构造子，可以视为等价于可数多个形如 $W\to W$ 的构造子。
（当然，这种无限性现在是类型理论\emph{内部}的，但对任何基础系统而言都理应如此。）
类似地，如果我们想要一个取“无限多个参数”的构造子，可以允许它取一个以某个无限类型参数化的参数族，例如 $(\nat\to W) \to W$，它取一个 $W$ 的元素的无限序列\index{sequence}。

\index{recursion principle!for an inductive type}%
那么，一旦有了这样一个归纳定义，可以用它来做什么呢？
首先，有一个\define{递归原理（recursion principle）}，它表明：为了定义一个函数 $f:W\to P$，只需考虑输入 $w:W$ 来自某个构造子的情形，并可在该构造子的输入上递归调用 $f$。
对于示例构造子~\eqref{eq:example-constructor}，我们要求 $P$ 配备一个类型为
\begin{narrowmultline}\label{eq:example-rechyp}
  d : (A\to W) \to (A\to P) \to (B\to C\to W) \to
  \narrowbreak
  (B\to C \to P) \to D \to W \to P \to P.
\end{narrowmultline}
的函数。
在这些假设下，递归原理给出 $f:W\to P$，并且它以显然的方式“保持构造子数据”——这就是计算规则，其中我们利用了输入的协变性\index{covariance}。
\index{computation rule!for inductive types}%
例如，在例子~\eqref{eq:example-constructor}中，计算规则说的是，对任意 $\alpha:A\to W$，$\beta:B\to C\to W$，$\delta:D$，以及 $\omega:W$，我们有
\begin{equation}
  f(c(\alpha,\beta,\delta,\omega)) \jdeq d(\alpha,f\circ \alpha,\beta, f\circ \beta, \delta, \omega,f(\omega)).\label{eq:example-comp}
\end{equation}
% As we have before in particular cases, when defining a particular function $f$, we may write these rules with $\defeq$ as a way of specifying the data $d$, and say that $f$ is defined by them.

\index{induction principle!for an inductive type}%
一般归纳类型 $W$ 的\define{归纳原理（induction principle）}只稍微复杂一点。
当然，我们从类型族 $P:W\to\type$ 出发，要求它配备“位于”$W$ 的构造子数据“之上”的构造子数据。
这意味着，上述诸如 $A\to P$ 这样的“递归调用”参数必须被替换为类型如 $\prd{a:A} P(\alpha(a))$ 的依值函数。
在~\eqref{eq:example-constructor} 的完整例子中，归纳原理对应的假设将要求
\begin{multline}\label{eq:example-indhyp}
d : \prd{\alpha:A\to W}\Parens{\prd{a:A} P(\alpha(a))} \to \narrowbreak
\prd{\beta:B\to C\to W} \Parens{\prd{b:B}{c:C} P(\beta(b,c))} \to\\
\prd{\delta:D}
\prd{\omega:W} P(\omega) \to
P(c(\alpha,\beta,\delta,\omega)).
\end{multline}
相应的计算规则形式上与~\eqref{eq:example-comp} 相同。
当然，递归原理是归纳原理当 $P$ 为常值族时的特例。
如前所述，归纳原理也称为\define{消去子（eliminator）}，而递归原理称为\define{非依值消去子（non-dependent eliminator）}。

如 \cref{sec:pattern-matching} 所述，我们也允许隐式地调用归纳与递归原理，为归纳原理的每个前提所对应的表达式写下带 $\defeq$ 的定义方程。
这称为通过（依值）\define{模式匹配（pattern matching）}来给出定义。
\index{pattern matching}%
\index{definition!by pattern matching}%
在我们这个贯穿始终的例子中，这意味着可以通过下式来定义 $f:\prd{w:W} P(w) $：
\[ f(c(\alpha,\beta,\delta,\omega)) \defeq \cdots \]
其中 $\alpha:A\to W$、$\beta:B\to C\to W$、$\delta:D$ 和 $\omega:W$ 是绑定在右侧表达式中的变量。
\index{variable}%
此外，右侧表达式可能涉及对 $f$ 的递归调用，形如 $f(\alpha(a))$、$f(\beta(b,c))$ 和 $f(\omega)$。
当用归纳原理重新表述这个定义时，我们将这类递归调用分别替换为新的变量 $\bar\alpha(a)$、$\bar\beta(b,c)$ 和 $\bar\omega$。
\begin{align*}
  \bar\alpha &: \prd{a:A} P(\alpha(a))\\
  \bar\beta &: \prd{b:B}{c:C} P(\beta(b,c))\\
  \bar\omega &: P(\omega).
\end{align*}
\symlabel{defn:induction-wtype}%
这样，我们可以写：
\[ f \defeq \ind{W}(P,\, \lam{\alpha}{\bar\alpha}{\beta}{\bar\beta}{\delta}{\omega}{\bar\omega} \cdots ) \]
其中 $\ind{W}$ 的第二个参数具有~\eqref{eq:example-indhyp} 的类型。

我们不准备尝试形式化地描述有效归纳定义的语法，及其所产生的归纳与递归原理和模式匹配规则。
这当然可以做到（事实上，如果要实现一个计算机证明助理，这是必须做的），但不能提供更多洞见。
经过一定练习后，便能自动为任何归纳定义推导出归纳与递归原理，乃至不假思索地运用它们。

\section{归纳类型的推广}
\label{sec:generalizations}

\index{type!inductive!generalizations}%
归纳类型的概念在类型论中已被研究多年，并有许多许多推广形式：归纳类型族、相互归纳类型、归纳-归纳（inductive-inductive）类型、归纳-递归（inductive-recursive）类型等等。
本节我们将概述其中的一部分，有几种将在本书后续内容中用到。
（在 \cref{cha:hits} 中，我们将更深入地研究一种非常不同的归纳类型推广，它是 \emph{同伦}类型论所特有的。）

这些推广大多涉及同时归纳定义多个类型。
一个我们已经见过的、非常简单的例子是余积 $A+B$。
如果每次想要考虑两个类型的余积时，都必须为 $\nat+\nat$、$\nat+\bool$、$\bool+\bool$ 等等分别写下独立的归纳定义，那将非常繁琐。
相反，我们做一个定义，其中 $A$ 和 $B$ 是代表类型的变量；
\index{variable!type}%
在类型论中，它们被称为\define{参数（parameters）}。%
\indexdef{parameter!of an inductive definition}
因此，严格地说，这样定义的结果不是一个单一类型，而是一个类型族 $+ : \type\to\type\to\type$，它取两个类型作为输入并产生它们的余积。
类似地，列表\index{type!of lists}类型 $\lst A$ 也是一个以类型 $A$ 为参数的类型族 $\lst{\blank}:\type\to\type$。

在数学中，这类事情显而易见，不值一提，但我们在这里提出来是为了与下一个例子作对比。
注意，每个类型 $A+B$ 都是\emph{各自独立地}归纳定义的，每个类型 $\lst A$ 也是如此。
\index{type!family of!inductive}%
\index{inductive!type family}%
相比之下，我们也可以考虑\emph{共同}归纳定义整个类型族 $B:A\to\type$。
区别在于，现在的构造子可能会改变指标 $a:A$，其结果是，我们不能说单个类型 $B(a)$ 是归纳定义的，只能说整个类型族是归纳定义的。

\index{type!of vectors}%
\index{vector}%
标准的例子是\emph{指定长度的列表}类型，传统上称为\define{向量（vectors）}。
我们固定一个参数类型 $A$，并对 $n:\nat$ 定义一个由以下构造子生成的类型族 $\vect n A$：

\begin{itemize}
\item 一个长度为零的向量 $\nil:\vect 0 A$，
\item 一个函数 $\cons:\prd{n:\nat} A\to \vect n A \to \vect{\suc (n)} A$。
\end{itemize}

与列表不同，向量（其元素取自固定类型 $A$）形成一个以长度为索引的类型族。
$A$ 是一个参数，而我们称 $n:\nat$ 是这个归纳族的\define{指标（index）}。
\indexdef{index of an inductive definition}%
单个类型如 $\vect3A$ 并不是归纳定义的：构造 $\vect3A$ 中元素的构造子从该族中的另一个类型取输入，例如 $\cons:A \to \vect2A \to \vect3A$。

\index{induction principle!for type of vectors}
\index{vector!induction principle for}
具体来说，归纳原理也必须涉及整个类型族；因此，前提与结论需对指标进行适当的量化。
在向量类型的情形下，其归纳原理是说，给定一个类型族 $C:\prd{n:\nat} \vect n A \to \type$，以及

\begin{itemize}
\item 一个元素 $c_\nil : C(0,\nil)$，和
\item 一个函数 \narrowequation{c_\cons : \prd{n:\nat}{a:A}{\ell:\vect n A} C(n,\ell) \to C(\suc(n),\cons(a,\ell))}，
\end{itemize}

则存在一个函数 $f:\prd{n:\nat}{\ell:\vect n A} C(n,\ell)$ 使得
\begin{align*}
  f(0,\nil) &\jdeq c_\nil\\
  f(\suc(n),\cons(a,\ell)) &\jdeq c_\cons(n,a,\ell,f(\ell)).
\end{align*}

\index{predicate!inductive}%
\index{inductive!predicate}%
归纳族的一个用途是归纳地定义\emph{谓词}。
例如，我们可以将谓词 $\mathsf{iseven}:\nat\to\type$ 定义为一个以 $\nat$ 为指标的归纳族，其构造子如下：

\begin{itemize}
\item 一个元素 $\mathsf{even}_0 : \mathsf{iseven}(0)$，
\item 一个函数 $\mathsf{even}_{ss} : \prd{n:\nat} \mathsf{iseven}(n) \to \mathsf{iseven}(\suc(\suc(n)))$。
\end{itemize}

换句话说，我们规定 $0$ 是偶数，并且如果 $n$ 是偶数，那么 $\suc(\suc(n))$ 也是偶数。
这些构造子``显然''没有提供任何办法来构造，比如说，$\mathsf{iseven}(1)$ 的一个元素；又因为 $\mathsf{iseven}$ 理应由这些构造子自由生成，所以这样的元素必然不存在。
（不过，严格证明 $\neg \mathsf{iseven}(1)$ 也并非完全平凡。）
$\mathsf{iseven}$ 的归纳原理是说，要证明关于所有偶数自然数的命题，只需对 $0$ 证明它，并验证在加二后仍然成立即可。

\index{mathematics!formalized}%
归纳定义的谓词在数学形式化与软件验证中被大量使用。
但本书中不会过多使用它们，仅在 \cref{sec:ordinals,sec:compactness-interval} 中有几处例外。

\index{type!mutual inductive}%
\index{mutual inductive type}%
另一个重要的特例是当归纳族的指标类型有限时。
此时，我们可以等价地将其表述为由\emph{相互归纳}定义的有限个类型。
例如，我们可以通过相互归纳来定义偶数自然数类型 $\mathsf{even}$ 和奇数自然数类型 $\mathsf{odd}$，其中 $\mathsf{even}$ 由构造子

\begin{itemize}
\item $0:\mathsf{even}$ 和
\item $\mathsf{esucc} : \mathsf{odd}\to\mathsf{even}$
\end{itemize}

生成，而 $\mathsf{odd}$ 则由一个构造子

\begin{itemize}
\item $\mathsf{osucc} : \mathsf{even}\to \mathsf{odd}$
\end{itemize}

生成。
注意 $\mathsf{even}$ 和 $\mathsf{odd}$ 是简单类型（而非类型族），但它们的构造子可以互相引用。
如果我们将这个定义表述为一个归纳类型族 $\mathsf{paritynat} : \bool \to \type$，其中 $\mathsf{paritynat}(\bfalse)$ 和 $\mathsf{paritynat}(\btrue)$ 分别代表 $\mathsf{even}$ 和 $\mathsf{odd}$，那么它的构造子会是：

\begin{itemize}
\item $0 : \mathsf{paritynat}(\bfalse)$，
\item $\mathsf{esucc} : \mathsf{paritynat}(\btrue) \to \mathsf{paritynat}(\bfalse)$，
\item $\mathsf{osucc} : \mathsf{paritynat}(\bfalse) \to \mathsf{paritynat}(\btrue)$。
\end{itemize}

当将其显式表述为相互归纳定义时，$\mathsf{even}$ 和 $\mathsf{odd}$ 的归纳原理是说，给定 $C:\mathsf{even}\to\type$ 和 $D:\mathsf{odd}\to\type$，以及

\begin{itemize}
\item $c_0 : C(0)$，
\item $c_s : \prd{n:\mathsf{odd}} D(n) \to C(\mathsf{esucc}(n))$，
\item $d_s : \prd{n:\mathsf{even}} C(n) \to D(\mathsf{osucc}(n))$，
\end{itemize}

则存在 $f:\prd{n:\mathsf{even}} C(n)$ 和 $g:\prd{n:\mathsf{odd}}D(n)$ 使得
\begin{align*}
  f(0) &\jdeq c_0\\
  f(\mathsf{esucc}(n)) &\jdeq c_s(g(n))\\
  g(\mathsf{osucc}(n)) &\jdeq d_s(f(n)).
\end{align*}
特别地，正如我们只能对归纳族``整体''进行归纳一样，我们必须同时对 $\mathsf{even}$ 和 $\mathsf{odd}$ 进行归纳。
本书同样不会过多使用相互归纳定义。

\index{type!inductive-inductive}%
\index{inductive-inductive type}%
一种更为激进的进一步推广是，允许定义一个类型族 $B:A\to \type$，其中不仅是类型 $B(a)$，连类型 $A$ 本身也是同一个大的归纳定义的一部分。
换言之，我们不仅像在归纳族中那样，为 $B(a)$ 指定可以接受其他 $B(a')$ 作为输入的构造子；同时也为 $A$ 本身指定构造子，而这些构造子可以接受 $B(a)$ 作为输入。
这可以看作是一种归纳族，其指标与被索引的类型是同时归纳定义的；也可以看作是一种相互归纳定义，其中一个类型可以依赖于另一个。
更复杂的依赖结构也是可能的。
一般而言，这些就称为\define{归纳-归纳定义}。
在本书的大部分内容中，我们不会使用它们，但其高阶变体（参见 \cref{cha:hits}）将在 \cref{cha:real-numbers} 的几个实验性例子中出现。

\index{type!inductive-recursive}%
\index{inductive-recursive type}%
我们要提及的最后一个推广是\define{归纳-递归定义}，在这种定义中，我们同时归纳地定义一个类型以及该类型上的一个\emph{递归}函数。
具体来说，我们先固定一个已知类型 $P$，然后为归纳类型 $A$ 给出构造子，同时利用这些构造子为 $A$ 带来的递归原理定义一个函数 $f:A\to P$——特别之处在于，$A$ 的构造子也被允许引用 $f$ 的值。
我们目前还不知道如何从同伦的视角证立这类定义，因此本书将不会使用它们。

\index{type!inductive|)}%

\section{相等类型与恒等系统}
\label{sec:identity-systems}

\index{type!identity!as inductive}%
现在我们要指出，在同伦类型论中处于核心地位的\emph{相等类型}，同样也可以被认为是归纳定义的。
具体来说，它们是一种带指标的``归纳族''，即 \cref{sec:generalizations} 意义下的归纳族。
事实上，将相等类型描述为归纳族有\emph{两种}方式，由此便产生了 \cref{cha:typetheory} 中所述的那两种归纳原理：道路归纳与基于点的道路归纳。

在两种定义中，类型 $A$ 都是一个参数。
对于第一种定义，我们归纳地定义一个族 $=_A : A\to A\to \type$，它有两个属于 $A$ 的指标，并由以下构造子给出：

\begin{itemize}
\item 对任意 $a:A$，有一个元素 $\refl a : a=_A a$。
\end{itemize}

通过与其它归纳族作类比，我们可以从这个定义中提取出归纳原理。
该原理陈述如下：给定任意 \narrowequation{C:\prd{a,b:A} (a=_A b) \to \type,} 以及 $d:\prd{a:A} C(a,a,\refl{a})$，则存在 \narrowequation{f:\prd{a,b:A}{p:a=_A b} C(a,b,p)} 使得 $f(a,a,\refl a)\jdeq d(a)$。
这正是相等类型的道路归纳原理。

对于第二种定义，我们将一个元素 $a_0:A$ 与 $A:\type$ 一并视为参数，然后归纳地定义一族 $(a_0 =_A \blank):A\to \type$，该族仅有\emph{一个}属于 $A$ 的指标，并由以下构造子给出：

\begin{itemize}
\item 一个元素 $\refl{a_0} : a_0 =_A a_0$。
\end{itemize}

请注意，由于 $a_0:A$ 被固定为参数，构造子 $\refl{a_0}$ 在归纳定义中并非以函数形式出现，而仅仅是一个元素。
该定义对应的归纳原理是说：给定 $C:\prd{b:A} (a_0 =_A b) \to \type$ 以及一个元素 $d:C(a_0,\refl{a_0})$，则存在 $f:\prd{b:A}{p:a_0 =_A b} C(b,p)$ 使得 $f(a_0,\refl{a_0})\jdeq d$。
这正是相等类型的基于点的道路归纳原理。

将相等类型视为归纳类型的观点，历史上曾引发一些困惑，原因在于 \cref{sec:bool-nat} 中提到的一种直觉：归纳类型的所有元素都应当通过反复应用其构造子来得到。
对于像 \bool 和 \nat 这类普通的归纳类型，情况确实如此：我们在 \cref{thm:allbool-trueorfalse} 中看到，\bool 的每个元素要么是 $\bfalse$，要么是 $\btrue$；类似地，我们也能证明 \nat 的每个元素要么是 $0$，要么是某个后继。

然而，对于相等类型，这一点却\emph{不}成立：它只有一个构造子 $\refl{}$，但并非每条道路都等于常值道路。
更准确地说，我们仅使用相等类型的归纳原理（无论哪一种）都无法证明：$a=_A a$ 的每个居民都等于 $\refl a$。
为了真正展示一个反例，我们需要某种额外的原则，例如泛等公理——回想一下，在 \cref{thm:type-is-not-a-set} 中，我们正是利用泛等公理展示了一条特殊的道路 $\bool=_{\type}\bool$，它不等于 $\refl{\bool}$。

\index{free!generation of an inductive type}%
\index{generation!of a type, inductive|(}%
关键在于，正如对同伦初始代数的研究所证实的那样，一个归纳定义应当被视为由其构造子\emph{自由生成}的。
当然，一个自由生成的结构所包含的元素可以多于其生成元：例如，由两个符号 $x$ 和 $y$ 生成的自由群中，不仅包含 $x$ 和 $y$，还包含诸如 $xy$、$yx^{-1}y$ 以及 $x^3y^2x^{-2}yx$ 这样的单词。
一般而言，自由结构的元素不仅通过应用生成元得到，也通过应用环境结构中的运算而得到——比如，如果我们讨论的是自由群，那就是群运算。

就归纳类型而言，我们所讨论的是自由生成的\emph{类型}——那么，一个类型的结构所具有的``运算''指的是什么呢？

如果像传统类型论那样，将类型视作\emph{集合}，那么类型上并没有这样的运算，因此我们自然期望归纳类型中除了其构造子所产生的元素之外，不含其他元素。

在同伦类型论中，我们将类型视作\emph{空间}或 $\infty$-群胚，%
\index{.infinity-groupoid@$\infty$-groupoid}
此时\emph{道路}上便有了许多运算（拼接、取逆等等）——这将在 \cref{cha:hits} 中变得重要——但\emph{对象}（元素）上仍然没有运算。
因此，对我们而言，诸如``$\bool$ 的每个元素要么是 $\bfalse$ 要么是 $\btrue$''、``$\nat$ 的每个元素要么是 $0$ 要么是某个后继''这样的陈述依然成立。

然而，正如我们在 \cref{cha:basics} 中看到的那样，将类型视为 $\infty$-群胚也意味着要将函数视作函子，这包括了类型族 $B:A\to\type$。
于是，被视为归纳类型族的相等类型 $(a_0 =_A \blank)$，实际上是一个\emph{自由生成的函子} $A\to\type$。
具体地说，它是由一个元素 $\refl{a_0}: F(a_0)$ 自由生成的函子 $F:A\to\type$。
而函子确实在对象上有运算，即 $A$ 的态射（道路）的作用。

范畴论中的\emph{米田引理}\index{Yoneda!lemma}告诉我们，对于任意范畴 $A$ 和对象 $a_0$，由 $F(a_0)$ 中的一个元素自由生成的函子，正是可表函子 $\hom_A(a_0,\blank)$。
因此，相等类型 $(a_0 =_A \blank)$ 应当就是这个可表函子，而我们确实也是这样看待它的：$(a_0 =_A b)$ 就是 $A$ 中从 $a_0$ 到 $b$ 的态射（道路）所构成的空间。

\index{generation!of a type, inductive|)}

\mentalpause

将相等类型视为归纳族的理由之一，是为了能够应用\cref{sec:appetizer-univalence,sec:htpy-inductive}中的唯一性原理。
具体地说，我们可以通过在 $A\times A$ 上给出另一个满足相同归纳原理的类型族，从而在等价的意义下刻画类型 $A$ 的相等类型族。
这就引出了下面的定义和定理。

\indexsee{system, identity}{identity system}%
\index{identity!system!at a point|(defstyle}%

\begin{defn}\label{defn:identity-systems}
  设 $A$ 是一个类型，$a_0:A$ 是其元素。
  \begin{itemize}
  \item $(A,a_0)$ 上的一个 \define{带点谓词（pointed predicate）}
    \indexdef{predicate!pointed}%
    \indexdef{pointed!predicate}%
    是一个配备了元素 $r_0:R(a_0)$ 的类型族 $R:A\to\type$。
  \item 对于带点谓词 $(R,r_0)$ 与 $(S,s_0)$，一个映射族 $g:\prd{b:A} R(b) \to S(b)$ 称为 \define{带点的}，如果 $g(a_0, r_0)=s_0$。
    我们有
    \[ \mathsf{ppmap}(R,S) \defeq \sm{g:\prd{b:A} R(b) \to S(b)} (g(a_0, r_0)=s_0).\]
  \item 一个 \define{在 $a_0$ 处的恒等系统（identity system at $a_0$）}
    是一个满足如下条件的带点谓词 $(R,r_0)$：对于任意的类型族 $D:\prd{b:A} R(b) \to \type$ 与 $d:D(a_0,r_0)$，都存在一个函数 $f:\prd{b:A}{r:R(b)} D(b,r)$ 使得 $f(a_0,r_0)=d$。
\end{itemize}
\end{defn}

\begin{thm}\label{thm:identity-systems}
  对于 $(A,a_0)$ 上的带点谓词 $(R,r_0)$，以下条件逻辑等价。
  \begin{enumerate}
  \item $(R,r_0)$ 是一个在 $a_0$ 处的恒等系统。\label{item:identity-systems1}
  \item 对于任意的带点谓词 $(S,s_0)$，类型 $\mathsf{ppmap}(R,S)$ 是可缩的。\label{item:identity-systems2}
  \item 对于任意的 $b:A$，函数 $\transfib{R}{\blank}{r_0} : (a_0 =_A b) \to R(b)$ 是一个等价。\label{item:identity-systems3}
  \item 类型 $\sm{b:A} R(b)$ 是可缩的。\label{item:identity-systems4}
  \end{enumerate}
\end{thm}

注意，\ref{item:identity-systems1}$\Leftrightarrow$\ref{item:identity-systems2}$\Leftrightarrow$\ref{item:identity-systems3} 之间的等价是 \cref{lem:homotopy-induction-times-3} 的一个版本，它针对的是被视为在 $A$ 的一个元素上变化的归纳族的相等类型 $a_0 =_A \blank$。
当然，条件\ref{item:identity-systems2}--\ref{item:identity-systems4}都是命题，因此逻辑等价蕴含真正的等价。
（条件\ref{item:identity-systems1}也是一个命题，但我们在此不证明这一点。）
另外需要注意，与条件\ref{item:identity-systems1}--\ref{item:identity-systems3}不同，陈述\ref{item:identity-systems4}并未提及 $a_0$ 或 $r_0$。

\begin{proof}
  首先，假设~\ref{item:identity-systems1} 并令 $(S,s_0)$ 为一个带点谓词。
  定义 $D(b,r) \defeq S(b)$ 和 $d\defeq s_0: S(a_0) \jdeq D(a_0,r_0)$。
  既然 $R$ 是一个恒等系统，我们有 $f:\prd{b:A} R(b) \to S(b)$ 满足 $f(a_0,r_0) = s_0$；因此 $\mathsf{ppmap}(R,S)$ 是有实例的。
  现在假设 $(f,f_r),(g,g_r) : \mathsf{ppmap}(R,S)$，定义 $D(b,r) \defeq (f(b,r) = g(b,r))$，并令 $d \defeq f_r \ct \opp{g_r} : f(a_0,r_0) = s_0 = g(a_0,r_0)$。
  那么同样由于 $R$ 是恒等系统，我们有 $h:\prd{b:A}{r:R(b)} D(b,r)$ 使得 $h(a_0,r_0) = f_r \ct \opp{g_r}$ 成立。
  根据函数外延性以及 $\Sigma$ 类型和道路类型中道路的刻画，这些数据给出一个等式 $(f,f_r) = (g,g_r)$。
  因此，$\mathsf{ppmap}(R,S)$ 是一个有实例的命题，从而可缩；故~\ref{item:identity-systems2} 成立。

  现在假设~\ref{item:identity-systems2}，并定义 $S(b) \defeq (a_0=b)$ 满足 $s_0 \defeq \refl{a_0}:S(a_0)$。
  那么 $(S,s_0)$ 是一个带点谓词，且 $\lamu{b:B}{p:a_0=b} \transfib{R}{p}{r} : \prd{b:A} S(b) \to R(b)$ 是从 $S$ 到 $R$ 的带点的映射族。
  由假设，$\mathsf{ppmap}(R,S)$ 是可缩的，因而是有实例的，故也存在一个从 $R$ 到 $S$ 的带点的映射族。
  两个方向的复合分别是从 $R$ 到 $R$ 以及从 $S$ 到 $S$ 的带点的映射族，由于 $\mathsf{ppmap}(R,R)$ 和 $\mathsf{ppmap}(S,S)$ 都是可缩的，故它们等于各自的恒等映射族。
  于是~\ref{item:identity-systems3} 成立。

  现在假设~\ref{item:identity-systems3}，条件~\ref{item:identity-systems4} 可由 \cref{thm:contr-paths} 得出，这里用到了 $\Sigma$ 类型保持等价这一事实（即 \cref{thm:total-fiber-equiv} 的``若''方向）。

  最后，假定~\ref{item:identity-systems4}，并令 $D:\prd{b:A} R(b)\to  \type$ 和 $d:D(a_0,r_0)$。
  我们可以将 $D$ 等价地表达为一个族 $D':(\sm{b:A} R(b)) \to \type$。
  现在，因为 $\sm{b:A} R(b)$ 是可缩的，我们有
  \[p:\prd{u:\sm{b:A} R(b)} (a_0,r_0) = u. \]
  此外，由于可缩类型的道路类型仍是可缩的，我们有 $p((a_0,r_0)) = \refl{(a_0,r_0)}$。
  定义 $f(u) \defeq \transfib{D'}{p(u)}{d}$，由此得到 $f:\prd{u:\sm{b:A} R(b)} D'(u)$，或等价地 $f:\prd{b:A}{r:R(b)} D(b,r)$。
  最后，我们有
  \[f(a_0,r_0) \jdeq \transfib{D'}{p((a_0,r_0))}{d} = \transfib{D'}{\refl{(a_0,r_0)}}{d} = d.\]
  因此，~\ref{item:identity-systems1} 成立。
\end{proof}

\index{identity!system!at a point|)}%

对于被视为在 $A$ 的两个元素上变化的族的相等类型 $=_A$，我们可以推导出类似的结果。

\index{identity!system|(defstyle}%

\begin{defn}
  一个在类型 $A$ 上的 \define{恒等系统（identity system）}
  是一个族 $R:A\to A\to \type$ 配备一个函数 $r_0:\prd{a:A} R(a,a)$，使得对于任何类型族 $D:\prd{a,b:A} R(a,b) \to \type$ 和 $d:\prd{a:A} D(a,a,r_0(a))$，都存在一个函数 $f:\prd{a,b:A}{r:R(a,b)} D(a,b,r)$ 满足对所有 $a:A$ 有 $f(a,a,r_0(a))=d(a)$。
\end{defn}

\begin{thm}\label{thm:ML-identity-systems}
  对于配备 $r_0:\prd{a:A} R(a,a)$ 的 $R:A\to A\to\type$，以下条件逻辑等价。
  \begin{enumerate}
  \item $(R,r_0)$ 是 $A$ 上的一个恒等系统。\label{item:MLis1}
  \item 对于所有 $a_0:A$，带点谓词 $(R(a_0),r_0(a_0))$ 是 $a_0$ 处的一个恒等系统。\label{item:MLis2}
  \item 对于任何 $S:A\to A\to\type$ 和 $s_0:\prd{a:A} S(a,a)$，类型
    \[ \sm{g:\prd{a,b:A} R(a,b) \to S(a,b)} \prd{a:A} g(a,a,r_0(a)) = s_0(a) \]
    是可缩的。\label{item:MLis3}
  \item 对于任何 $a,b:A$，映射 $\transfib{R(a)}{\blank}{r_0(a)} : (a =_A b) \to R(a,b)$ 是一个等价。\label{item:MLis4}
  \item 对于任何 $a:A$，类型 $\sm{b:A} R(a,b)$ 是可缩的。\label{item:MLis5}
  \end{enumerate}
\end{thm}

\begin{proof}
  \ref{item:MLis1}$\Leftrightarrow$\ref{item:MLis2} 的证明完全类似于相等类型的道路归纳原理与基于点的道路归纳原理之间的等价性证明；参见 \cref{sec:identity-types}。
  与\ref{item:MLis4}和\ref{item:MLis5}的等价则由 \cref{thm:identity-systems} 得出，而\ref{item:MLis3}是直接的。
\end{proof}

\index{identity!system|)}%

这个刻画之所以有趣，原因之一是它提供了陈述泛等性和函数外延性的另一种方式。
\index{univalence axiom}%
泛等公理针对宇宙 \UU 所说的正是：类型族
\[ (\eqv{\blank}{\blank}) : \UU\to\UU\to\UU \]
配上 $\idfunc : \prd{A:\UU} (\eqv AA)$ 后满足~\cref{thm:ML-identity-systems}\ref{item:MLis4}。
因此，它等价于~\ref{item:MLis1}的对应形式，我们可以将其陈述如下。

\begin{cor}[等价归纳]\label{thm:equiv-induction}
  \index{induction principle!for equivalences}%
  \index{equivalence!induction}%
  给定任何类型族 \narrowequation{D:\prd{A,B:\UU} (\eqv AB) \to \type} 和函数 $d:\prd{A:\UU} D(A,A,\idfunc[A])$，存在 \narrowequation{f : \prd{A,B:\UU}{e:\eqv AB} D(A,B,e)} 使得对所有 $A:\UU$ 有 $f(A,A,\idfunc[A]) = d(A)$。
\end{cor}

换言之，要证明关于所有等价的性质，只需对恒等映射证明它即可。
我们已经在 \cref{lem:qinv-autohtpy} 中使用过这个原理（但未作一般性陈述）。

类似地，函数外延性说的是：对于任何 $B:A\to\type$，类型族
\[ (\blank\htpy\blank) : \Parens{\prd{a:A} B(a)} \to \Parens{\prd{a:A} B(a)} \to \type
\]
配上 $\lamu{f:\prd{a:A} B(a)}{a:A} \refl{f(a)}$ 后满足~\cref{thm:ML-identity-systems}\ref{item:MLis4}。
因此，它也等价于~\ref{item:MLis1}的对应形式。

\begin{cor}[同伦归纳]\label{thm:htpy-induction}
  \index{induction principle!for homotopies}%
  \index{homotopy!induction}%
  给定任何 \narrowequation{D:\prd{f,g:\prd{a:A} B(a)} (f\htpy g) \to \type} 和 $d:\prd{f:\prd{a:A} B(a)} D(f,f,\lam{x}\refl{f(x)})$，存在
  %
  \begin{equation*}
    k:\prd{f,g:\prd{a:A} B(a)}{h:f\htpy g} D(f,g,h)
  \end{equation*}
  %
  使得对所有 $f$ 有 $k(f,f,\lam{x}\refl{f(x)}) = d(f)$。
\end{cor}

\sectionNotes

归纳定义在数学中历史悠久，可以说至少可追溯至 Frege（弗雷格）和 Peano（皮亚诺）的自然数公理。\index{Frege}\index{Peano} %
更为一般的“归纳谓词”也并不罕见，但在集合论基础中，它们通常被显式构造出来——或是作为某个适当的子集类的交，或是沿着序数进行超限迭代——而非被视为一个基本的概念。

在类型论中，归纳定义的特例可追溯至 Martin-L\"of（马丁-洛夫）的原始论文：\cite{martin-lof-hauptsatz} 提出了一种归纳定义的谓词与关系的一般性概念；归纳类型的概念在 Martin-L\"of 最早的类型论论文 \cite{Martin-Lof-1973} 中已经出现（但仅有具体实例，并未作为一般性概念提出）；
随后在 \cite{Martin-Lof-1979} 中，才作为与 $\w$-类型相关的一般性概念出现。\index{Martin-L\"of}%

归纳类型的一般概念由 Constable 与 Mendler 于 1985 年引入 \cite{DBLP:conf/lop/ConstableM85}。内涵类型论中归纳类型的一般图式则由
\cite{PfenningPaulinMohring, CoquandPaulin, Dybjer:1991} 提出。

归纳-递归定义的概念出现于 \cite{Dybjer:2000}。一个重要的类型论概念是树类型（这是 Post 系统概念在类型论中的一般表述），它出现于 \cite{PeterssonSynek}。

自然数作为 $\nat$-代数范畴中始对象的\emph{泛性质}，归功于 Lawvere（劳维尔）\cite{lawvere:adjinfound}\index{Lawvere}。
后来，\cite{mp:wftrees} 将此推广，把 $\w$-类型描述为多项式自函子的初始代数。\index{endofunctor!algebra for}
此类泛性质与相应归纳原理之间的相干同伦等价性（\cref{sec:initial-alg,sec:htpy-inductive}）归功于 \cite{ags:it-hott}。

关于类型论的同伦论语义中归纳类型的具体构造，可参见 \cite{klv:ssetmodel,mvdb:wtypes,ls:hits}。

\sectionExercises

\begin{ex}\label{ex:ind-lst}
  根据 $\lst{A}$ 在 \cref{sec:bool-nat} 中作为归纳类型的定义，推导其归纳原理。\index{type!of lists}
\end{ex}

\begin{ex}\label{ex:same-recurrence-not-defeq}
  构造自然数上的两个函数，使它们在判断相等意义下满足同一递归式\index{recurrence} $(e_z, e_s)$，但彼此并不判断相等。
\end{ex}

\begin{ex}\label{ex:one-function-two-recurrences}
  在同一类型 $E$ 上构造两个不同的递归式 $(e_z,e_s)$，使得同一函数 $f:\nat\to E$ 在判断相等意义下同时满足二者。
\end{ex}

\begin{ex}\label{ex:bool}
  证明：对任意类型族 $E : \bool \to \type$，归纳算子
  \[ \ind{\bool}(E) : \big(E(\bfalse) \times E(\btrue)\big) \to \prd{b : \bool} E(b) \]
  都是一个等价。
\end{ex}

\begin{ex}\label{ex:ind-nat-not-equiv}
  证明：对于 $\nat$ 而言，与 \cref{ex:bool} 类似的命题不成立。
\end{ex}

\begin{ex}\label{ex:no-dep-uniqueness-failure}
  证明：如果我们对 $\w$-类型假设的是简单消去（而非依值消去），则唯一性性质（\cref{thm:w-uniq} 的类比）不再成立。
  亦即，举出一个满足 $\w$-类型递归原理的类型，但其中的函数并不由其递归式\index{recurrence}唯一确定。
\end{ex}

\begin{ex}\label{ex:loop}
  假设在 \cref{sec:strictly-positive} 开头处类型 $C$ 的“归纳定义”中，我们将类型 \nat 替换为 \emptyt。
  类似于 \eqref{eq:fake-recursor}，我们可以为此类型考虑一个递归原理，其假设为
  \[ h:(C\to\emptyt) \to (P\to\emptyt) \to P. \]
  试证：即使没有计算规则，这个递归原理也是不一致的，即它允许我们构造出 \emptyt 中的一个元素。
\end{ex}

\begin{ex}\label{ex:loop2}
  现在考虑一个“归纳类型” $D$，它只有一个构造子 $\mathsf{scott}:(D\to D) \to D$。
  \index{Scott}%
  \cref{sec:strictly-positive} 中为 $C$ 提议的第二个递归子导出 $D$ 的如下递归子：
  \[ \rec{D} : \prd{P:\UU} ((D\to D) \to (D\to P)\to P) \to D \to P \]
  其计算规则为 $\rec{D}(P,h,\mathsf{scott}(\alpha)) \jdeq h(\alpha,(\lam{d} \rec{D}(P,h,\alpha(d))))$。
  试证这同样会导致矛盾。
\end{ex}

\begin{ex}\label{ex:inductive-lawvere}
  设 $A$ 为任意类型，并考虑类型 $L_A$ 的一个一般性“归纳定义”，其构造子为 $\mathsf{lawvere}:(L_A\to A) \to L_A$。
  \cref{sec:strictly-positive} 中为 $C$ 提议的第二个递归子导出 $L_A$ 的如下递归子：
  \[ \rec{L_A} : \prd{P:\UU} ((L_A\to A) \to P) \to L_A\to P \]
  其计算规则为 $\rec{L_A}(P,h,\mathsf{lawvere}(\alpha)) \jdeq h(\alpha)$。
  利用此递归子，证明 $A$ 具有\define{不动点性质（fixed-point property）}，即对任意函数 $f:A\to A$，均存在 $a:A$ 使得 $f(a)=a$。
  \index{Lawvere}%
  \index{fixed-point property}%
  特别地，若 $A$ 是一个不具有不动点性质的类型（例如 $\emptyt$、$\bool$ 或 $\nat$），则 $L_A$ 是不一致的。
\end{ex}

\begin{ex}\label{ex:ilunit}
  承接 \cref{ex:inductive-lawvere}，考虑 $L_{\unit}$；由于 $\unit$ 确实具有不动点性质，$L_{\unit}$ 并非显然不一致。
  试为 $L_{\unit}$ 表述一个归纳原理及其计算规则（类似于其递归子），并利用此归纳原理证明 $L_{\unit}$ 是可缩的。
\end{ex}

\begin{ex}\label{ex:empty-inductive-type}
在 \cref{sec:bool-nat} 中，我们定义了某个类型 $A$ 中元素的有限列表类型 $\lst A$。
现考虑类似的归纳定义，定义类型 $\lost A$，其唯一的构造子为
\[ \cons: A \to \lost A \to \lost A. \]
证明 $\lost A$ 等价于 $\emptyt$。
\end{ex}

\begin{ex}\label{ex:Wprop}
  设 $A$ 是一个命题，且 $B:A\to \UU$。
  \begin{enumerate}
  \item 证明 $\wtype{a:A} B(a)$ 是一个命题。
  \item 证明 $\wtype{a:A} B(a)$ 等价于 $\sm{a:A} \neg B(a)$。
  \item 不使用 $\wtype{a:A} B(a)$，证明 $\sm{a:A} \neg B(a)$ 是 \cref{sec:htpy-inductive} 意义下的同伦 $\w$-类型 $\wtypeh{a:A} B(a)$。
  \end{enumerate}
\end{ex}

\begin{ex}\label{ex:Wbounds}
  设 $A:\UU$ 且 $B:A\to \UU$。
  \begin{enumerate}
  \item 证明 $\Parens{\sm{a:A} \neg B(a)} \to \Parens{\wtype{a:A} B(a)}$。
  \item 证明 $\Parens{\wtype{a:A} B(a)} \to \Parens{\neg \prd{a:A} B(a)}$。
  \end{enumerate}
\end{ex}

\begin{ex}\label{ex:Wdec}
  设 $A:\UU$，并假定 $B:A\to \UU$ 是可判定的，即 $\prd{a:A} (B(a)+\neg B(a))$（见 \cref{defn:decidable-equality}）。
  证明 $\Parens{\wtype{a:A} B(a)} \to \Parens{\sm{a:A} \neg B(a)}$。
\end{ex}

\begin{ex}\label{ex:Wbounds-loose}
  证明下列命题是逻辑等价的。
  \begin{enumerate}
  \item 对任意 $A:\set$ 和 $B:A\to \prop$，有 $\Parens{\wtype{a:A} B(a)} \to \Brck{\sm{a:A} \neg B(a)}$。\label{item:Wbounds-loose-Sigma}
  \item 对任意 $A:\set$ 和 $B:A\to \prop$，有 $\Parens{\neg \prd{a:A} B(a)} \to \Brck{\wtype{a:A} B(a)}$。\label{item:Wbounds-loose-Pi}
  \item 排中律成立（如 \cref{sec:intuitionism} 中所述）。
  \end{enumerate}
  类似地，利用 \cref{thm:not-lem}，证明：若假设 \ref{item:Wbounds-loose-Sigma} 或 \ref{item:Wbounds-loose-Pi} 中的蕴含式对所有 $A:\UU$ 和 $B:A\to \UU$ 成立，将导致不一致。
\end{ex}

\begin{ex}\label{ex:Wimpred}
  对于 $A:\UU$ 和 $B:A\to \UU$，定义
  \[ W'_{A,B} \defeq \prd{R:\UU} \Parens{\prd{a:A} (B(a) \to R) \to R} \to R \]
  $W'_{A,B}$ 称为 $\wtype{a:A} B(a)$ 的\define{非直谓编码（impredicative encoding）}。
  \index{impredicative!encoding of a W-type@encoding of a $\w$-type}%
  \index{W-type@$\w$-type!impredicative encoding of}%
  注意，与 $\wtype{a:A} B(a)$ 不同，它位于比 $A$ 和 $B$ 更高的宇宙中。
  \begin{enumerate}
  \item 证明 $W'_{A,B}$ 逻辑等价于 $\wtype{a:A} B(a)$（逻辑等价的定义见 \cref{sec:pat}）。
  \item 证明 $W'_{A,B}$ 蕴含 $\neg\neg \sm{a:A} \neg B(a)$。
  \item 不使用 $\wtype{a:A} B(a)$，证明 $W'_{A,B}$ 满足与 $\wtype{a:A} B(a)$ 相同的\emph{递归}原理，可用于定义取值于宇宙 $\UU$（$W'_{A,B}$ 自身不属于该宇宙）中类型的函数。
  \item 利用\LEM，给出一个 $A:\UU$ 和 $B:A\to \UU$ 的例子，使得 $W'_{A,B}$ 不与 $\wtype{a:A} B(a)$ 等价。
  \end{enumerate}
\end{ex}

\begin{ex}\label{ex:no-nullary-constructor}
  证明对于任意 $A:\UU$ 和 $B:A\to\UU$，有
  \[ \eqv{\neg\Parens{\wtype{a:A} B(a)}}{\neg\Parens{\sm{a:A} \neg B(a)}}. \]
  换言之，$\wtype{a:A} B(a)$ 为空当且仅当它没有零元构造子。
  （与 \cref{ex:empty-inductive-type} 比较。）
\end{ex}

% Local Variables:
% TeX-master: "hott-online"
% End: